% Source - https://tex.stackexchange.com/a/630315
% Posted by Mico, modified by community. See post 'Timeline' for change history
% Retrieved 2026-10-02, License - CC BY-SA 4.0

\documentclass{amsart}
\usepackage[T1]{fontenc}
\usepackage{amssymb}
\newtheorem{theorem}[subsection]{Theorem}

\title[RH Euclidean Cartesian]{Proving the Riemann Hypothesis \\
	with Euclidean Distance, Cartesian Coordinates and Series Divergence}
\author{Zinoo Park  \\ \MakeLowercase{zinoopark09@alumni.harvard.edu}}
\thanks{Property of BigPlace LLC, Wyoming, USA}
\date{20261013}

\begin{document}
	\maketitle
	
	\begin{abstract}
		The Euclidean distance and Cartesian coordinates are used to show that the real part of a non-trivial zero, $Re(s)$, of the Riemann zeta function is always $\frac12$. If the distance between two points on a 2D plane equals zero, they occupy the same coordinates. Let $z=\zeta(s)=X+Yi$ where $s=\frac{1}{b}+ci$. Through analytic continuation where $D_{s}=1-2^{1-s}=v+wi$ and $D_{s}\zeta(s) = \eta(s)$, we transform $z$ into $Z = (vX-wY) + (wX+vY)i = \eta(s)$. On an XY plane, let $Z=(vX-wY, wX+vY)$ and the target point $\Omega=(p, q)$. If $Z\Omega = 0$ when $\Omega=(0,0)$, it is implied that $\eta(s)=0$. Rearranging the terms leads to $\zeta(\frac{2}{b})$, equivalent to the negative sum of all cross terms in $X^2 + Y^2$ that diverges given $0<\frac{1}{b}<1$. Further analytic continuation yields $\zeta(\frac{2}{b}) = \frac{\eta(2/b)}{D_{2/b}}$ equalling a similar minus sum of cross terms from $G^2 + H^2$ where $\eta(s)= G + Hi$. The latter diverges while $b \geq 2$. Under the initial and additional conditions, $\frac{\eta(2/b)}{D_{2/b}}$ diverges only when $D_{2/b} = 1-2^{1-2/b} = 0$ leading to $b=2$.
	\end{abstract}
	
	\section{Introduction}
	Given any $s$ where $s \in \mathbb{C}$ and $0<Re(s)<1$, the series yielded by the Riemann zeta function does not converge unless via analytic continuation. This means that if $\zeta(s) = z = X+Yi$, $X$ and $Y$ do not have fixed numerical values. In such a case, whether or not $z(X,Y)$ and another point $\Omega(p,q)$, where $p,q \in \mathbb{R}$, occupy the same coordinates on an XY plane cannot be determined by simply equating one to the other. The equation when doing so produces false statements as neither $X-p$ nor $Y-q$ can equate to zero, let alone any other real number.
	
	\[ z(X,Y) = \Omega(p,q) \quad \Rightarrow \quad X=p, Y=q \]
	% \[ X-p \neq 0, Y -q \neq 0 \qquad\qquad (\because \, X = \text{DNE}, \quad Y = \text{DNE}) \]
	\[ X-p \neq 0, Y -q \neq 0 \qquad\qquad (\because \, X, Y \quad \text{DNE}) \]
	\[ \Rightarrow z(X,Y) \neq \Omega(p,q) \]
	
	There is another method of checking for overlap in coordinates. It is to calculate the Euclidean distance between the two given points and compare it to zero.
	
	\[ z\Omega = \sqrt{(X-p)^2 + (Y-q)^2} = 0 \]
	
	The terms before expansion is the reason we simply equate one pair of coordinates to another. If all terms involved were real numbers, the squares of differences, namely $(X-p)^2$ and $(Y-q)^2$, would individually be equal to or above zero. The only way the result becomes zero is when both of the differences are zero.
	
	\[ X-p=0 \qquad \Rightarrow \qquad X=p  \]
	\[ Y-q=0 \qquad \Rightarrow \qquad Y=q  \]
	
	However, as we have already established that this is not the case, we must consider other possibilities. Even when some series diverge on their own, when combined through addition or subtraction, they may converge. Moreover, the self-Cauchy products of divergent series may yield sequences that cancel out others when together, leading to convergence. For this, we expand $z\Omega$.
	
	\[ z\Omega = \sqrt{X^2 - 2pX + p^2 + Y^2 - 2qY + q^2} = 0 \]
	
	One glaring issue in this expanded version is the divergence of $X$ and $Y$. We may suppose that the two Cauchy products, $X^2$ and $Y^2$, converge as a group $X^2+Y^2$, but we still cannot guarantee that $-2pX-2qY$ converge by cancelling each other out. The easiest option that anyone would be tempted to do would be to let $p=0$ and $q=0$.
	
	\[ \text{If} \quad p=0, q=0 \qquad \Rightarrow \qquad \Omega(0,0) \]
	
	\[ z\Omega = \sqrt{X^2 - 2 \cdot 0 \cdot X + 0^2 + Y^2 - 2 \cdot 0 \cdot Y + 0^2} = 0 \]
	
	\[ = \sqrt{X^2 + Y^2} = 0 \]
	
	This implies that we are making $\Omega$ origin $O(0,0)$. Subsequently, this means that $z(X,Y)$ equals or converges to the origin, a paradox when considering what was initially stated for X and Y. This is where analytic continuation steps in when the concept discussed in this paper is applied to the Riemann zeta function and its non-trivial zeros.
	
	In the sections that follow, by computing the Euclidean distance in a 2D Cartesian coordinate system, we will observe the characteristics of resulting terms. This will help us understand values from $\zeta(s)$ where $s = \frac{1}{b} + ci$ and properties of $\zeta(\frac{2}{b})$.
	\[\]
	
	\section{Definition}
	\subsection{The Riemann zeta and Dirichlet eta}
	We may define the functions Riemann zeta and the Dirichlet eta as below where $ k \in \mathbb{N} $.
	\[ \zeta(s) := \frac{1}{1^s} + \frac{1}{2^s} + \frac{1}{3^s} + \frac{1}{4^s} + {\dots} + \frac{1}{k^s} = \sum\limits_{k=1}^{\infty}\frac{1}{k^s} \]
	\[ \eta(s) := \frac{1}{1^s} - \frac{1}{2^s} + \frac{1}{3^s} - \frac{1}{4^s} + ... \pm \frac{1}{k^s} = \sum\limits_{k=1}^{\infty}\frac{(-1)^{k-1}}{k^s} \qquad k \in \mathbb{N} \]
	Let $s$ be a complex number within the critical strip, a non-trivial zero of the Riemann zeta function expressed in real numbers $b$ and $c$.
	\[ s \in \mathbb{C}, \qquad s = \sigma + ci = \frac{1}{b} + ci \]
	\[ b,c \in \mathbb{R}, \qquad c \neq 0 \]
	\[ 0 < \sigma < 1 \quad \Rightarrow \quad 0 < \frac{1}{b} < 1 \quad \text{or} \quad 1<b<\infty \]
	\[\]	
	
	\subsection{The denominator $D$ for analytic continuation given s}
	The term $D_s$ enables analytic continuation between the $\zeta(s)$ and $\eta(s)$, and may be expressed as a complex number. It is important to know how its real part $v$ and imaginary part $w$ behave as we continue having to arrange various equations including them.
	\[ \zeta(s) = \frac{\eta(s)}{D_s} \quad \text{or} \quad D_s\zeta(s) = \eta(s) \]
	\[ D_{s} = 1-2^{1-s} = v + wi \]
	\[ \bar{D}_{s} = v - wi,  \qquad v,w \in \mathbb{R} \]
	\[\]
	
	\subsubsection{Range of $v^2+w^2$}
	Let $1 - \sigma = t$ to observe the properties of real numbers $v$ and $w$.
	\[ 1 - \sigma = t \]
	\[ D_{s} = 1-2^{1-s} \]
	\[ 2^{1-s} = 2^{1-\sigma - ci} = 2^{t-ci} = 2^t \cdot 2^{-ci} \]
	\[ 2^{-ci} = e^{\ln2 \cdot (-c)i} = e^{(-c \ln 2) i} = \cos(-c \ln 2) + \sin(-c \ln 2)i \]
	\[ 2^t \cdot 2^{-ci} = 2^t(\cos(-c \ln 2) + \sin(-c \ln 2)i) \]
	\[ D_{s} = 1-2^t \cdot 2^{-ci} = 1 - 2^t\cos(-c \ln 2) - 2^t\sin(-c \ln 2)i \]
	\[ = 1 - 2^t\cos(c \ln 2) + 2^t\sin(c \ln 2)i \]
	\[ = v + wi \]
	\[ v = 1 - 2^t\cos(c \ln 2) = 1 - 2^{1-\sigma}\cos(c \ln 2) \]
	\[ w = 2^t\sin(c \ln 2) = 2^{1-\sigma}\sin(c \ln 2) \]
	\[\]
	
	Given $ 0 < \sigma < 1 $, ranges of $1-\sigma$ and $2^{1-\sigma}$ can be determined.
	\[ 0 < 1-\sigma < 1 \qquad \Rightarrow \qquad 1 < 2^{1-\sigma} < 2 \]
	
	Suppose $v=0$ so that we may know what $w$ cannot equal.
	\[ v = 1 - 2^{1-\sigma}\cos(c \ln 2) = 0 \]
	\[ 1 = 2^{1-\sigma}\cos(c \ln 2) \]
	\[ \frac{1}{2} < \cos(c \ln 2) < 1 \quad (\because 1 < 2^{1-\sigma} < 2) \]
	
	Since $\cos^2\theta + \sin^2\theta = 1$ or $ \sin\theta = \pm \sqrt{1-\cos^2{\theta}} $,
	\[ \frac{1}{4} < \cos^2(c \ln 2) < 1 \quad \Rightarrow \quad 0 < 1 - \cos^2(c \ln 2) < \frac{3}{4} \]
	\[ 1 - \cos^2(c \ln 2) \neq 0 \quad \Rightarrow \quad \pm\sqrt{1 - \cos^2(c \ln 2)} = \sin(c \ln 2) \neq 0 \]
	\[ \Rightarrow w = 2^{1-\sigma}\sin(c \ln 2) \]
	\[ \therefore w \neq 0 \quad (\because 2^{1-\sigma} \neq 0) \]
	
	If $w=0$,
	\[ w = 2^{1-\sigma}\sin(c \ln 2) = 0 \]
	\[ \sin(c \ln 2) = 0 \qquad (\because 2^{1-\sigma} \neq 0) \]
	\[ c \ln 2 = l\cdot\pi \qquad (l \in \mathbb{Z}) \]
	\[ \Rightarrow \cos(l\pi) = \pm 1 \]
	\[ \Rightarrow v = 1 - 2^{1-\sigma} \cdot (1) = 1 \quad \text{or} \quad v = 1 - 2^{1-\sigma} \cdot (-1) \]	
	\[ \Rightarrow -1 < v < 0 \quad \text{or} \quad 2 < v < 3  \]	
	\[ \therefore v \neq 0 \]
	\[\]
	Therefore, we can deduce that $v^2 + w^2 \neq 0$ and that $v^2 + w^2 > 0$.
	\[\]
	
	\subsubsection{$D_s$ as a limit function of s}
	Despite the fact $D_{s} = 1-2^{1-s}$, it may be more than a complex number. It is a limit function of $s$. The terms $1-2^{1-s}$ are derived from the relationship between functions Riemann zeta and Dirichlet eta. Regardless of convergence or divergence, the concept or usage of a limit function is implied in each series. Therefore, even when $D_s$ can be expressed as $1-2^{1-s}$ or $v+wi$, its inherent property of being one should be remembered and utilised.
	
	\[ \zeta(s) = \lim_{n \to \infty} \zeta_{n}(s) = \lim_{n \to \infty} \sum\limits_{k=1}^{n}\frac{1}{k^s}  \]
	\[ \eta(s) = \lim_{n \to \infty} \eta_{n}(s) = \lim_{n \to \infty} \sum\limits_{k=1}^{n}\frac{(-1)^{k-1}}{k^s}  \]
	
	\[ D_s = \frac{\eta(s)}{\zeta(s)} = \frac{\lim\limits_{n \to \infty} \eta_{n}(s)}{\lim\limits_{n \to \infty} \zeta_{n}(s)} = \frac{\lim\limits_{n \to \infty} \sum\limits_{k=1}^{n}\frac{(-1)^{k-1}}{k^s}}{\lim\limits_{n \to \infty} \sum\limits_{k=1}^{n} \frac{1}{k^s}} \]
	\[= \lim\limits_{n \to \infty} \frac{\sum\limits_{k=1}^{n}\frac{(-1)^{k-1}}{k^s}}{\sum\limits_{k=1}^{n} \frac{1}{k^s}} = \lim\limits_{n \to \infty}D(s) \]
	\[\]
	
	\subsection{Expressing $\zeta(s)$ and its summands}
	Each term of the series from $\zeta(s)$ can be expressed in $b$ and $c$.
	\[  \frac{1}{k^s} = \frac{1}{k^{\frac{1}{b} + ci}} = \frac{\cos({c\ln k})}{\sqrt[b]{k}} + \frac{-\sin({c\ln k})}{\sqrt[b]{k}} i \]
	Then, $\zeta(s)$ can be represented as a complex number $z = X + Yi$. It is important to note that $z$ may also be a point on an XY plane.
	\[ \zeta(s) \qquad = \qquad \frac{1}{1^s} \qquad \Big|  \qquad \frac{\cos({c\ln 1})}{\sqrt[b]{1}} \qquad \frac{-\sin({c\ln 1})}{\sqrt[b]{1}} \quad i  \]
	
	\[ \quad \qquad  \qquad \qquad \frac{1}{2^s} \qquad \Big|  \qquad \frac{\cos({c\ln 2})}{\sqrt[b]{2}} \qquad \frac{-\sin({c\ln 2})}{\sqrt[b]{2}} \quad i  \]
	
	\[ \quad \qquad  \qquad \qquad \frac{1}{3^s} \qquad \Big|  \qquad \frac{\cos({c\ln 3})}{\sqrt[b]{3}} \qquad \frac{-\sin({c\ln 3})}{\sqrt[b]{3}} \quad i  \]
	
	\[ \vdots \]
	
	\[ \qquad \qquad \quad + \quad \frac{1}{k^s} \qquad \Big|  \qquad \frac{\cos({c\ln k})}{\sqrt[b]{k}} \qquad \frac{-\sin({c\ln k})}{\sqrt[b]{k}} \quad i  \]
	
	\[\quad \rule{12 cm}{.8pt} \quad \]
	\[ \qquad \qquad \qquad z \qquad\qquad =  \qquad\qquad X \qquad + \qquad Y \quad i \qquad\qquad \]
	\[ \text{or} \]
	\[ z (X, Y) \]
	\[\]
	Where $k \to \infty$, $X$ and $Y$ are the sums of $x_k$ and $y_k$, respectively, and every term including a cosine value becomes $x_k$ whereas every term with a sine value becomes $y_k$. 
	\[ X = \sum\limits_{k=1}^{\infty}x^k = Re(\zeta(s)) \, , \quad Y = \sum\limits_{k=1}^{\infty}y^k = Im(\zeta(s)) \]
	\[\]
	\[ x_k = \frac{\cos({c\ln k})}{\sqrt[b]{k}} = \frac{\cos(-{c\ln k})}{\sqrt[b]{k}} = \frac{\cos({c\ln k^{-1}})}{\sqrt[b]{k}} = \frac{\cos({c\ln \frac{1}{k}})}{\sqrt[b]{k}} \]
	\[ y_k = \frac{-\sin({c\ln k})}{\sqrt[b]{k}} = \frac{\sin({-c\ln k})}{\sqrt[b]{k}} = \frac{\sin({c\ln k^{-1}})}{\sqrt[b]{k}} = \frac{\sin({c\ln \frac1k })}{\sqrt[b]{k}} \]
	
	
	
	
	
	
	
	\section{Calculating the Euclidean distance}
	\subsection{Applying analytic continuation}
	Since $\eta(s) = D_s\zeta(s)$, let $D_s = 1-2^{1-s}=v+wi$ and multiply it by $\zeta(s)$ in order to obtain $\eta(s)$ in terms of $X$ and $Y$. Then let $Z = \eta(s)$, the newly transformed complex value.
	\[ D_s\zeta(s) = \eta(s) = (v+wi) \times (X+Yi) = vX + vYi + wXi - wY  \]
	\[ = (vX - wY) + (wX + vY)i = Z \]
	\[ Z = (vX - wY) + (wX + vY)i \]
	\[ \text{or} \]
	\[ Z \Bigl( (vX - wY), (wX + vY) \Bigr) \]
	
	\subsection{The Euclidean distance}
	Let the target point $\Omega$ be a point with coordinates $(p,q)$ in a 2D Cartesian coordinate system.
	\[ \Omega(p, q) \ , \quad p,q \in \mathbb{R} \]
	
	The Euclidean distance between $Z$ and $\Omega$ on the given 2D plane will be $Z\Omega$.
	\[ Z\Omega = \sqrt{\{(vX-wY)-p\}^2 + \{(wX+vY)-q\}^2} \]
	If the distance is equal to zero, implying that the two points occupy the same coordinates, this can be expanded as follows.
	\[ Z\Omega = \sqrt{\{(vX-wY)-p\}^2 + \{(wX+vY)-q\}^2} = 0 \]
	\[ = \sqrt{\{ (vX-wY)^2 {+} p^2 -2p(vX-wY) \} {+} \{ (wX{+}vY)^2 {+} q^2 {+} 2q(wX{+}vY) \}} = 0 \]
	We can also imply that $\eta(s)=Z$ equals or converges to zero by moving point $\Omega$ to the origin$(0,0)$ if we let $p=0$ and $q=0$.
	\[ Z\Omega = \sqrt{\{ (vX-wY)^2 {+} 0^2 -2 \cdot 0 (vX-wY) \} {+} \{ (wX{+}vY)^2 {+} 0^2 {+} 2 \cdot 0 (wX{+}vY) \}} \]
	\[ = \sqrt{(vX-wY)^2 {+} (wX{+}vY)^2} = 0 \]
	
	
	When analytic continuation is applied and $Z = (vX - wY) + (wX + vY)i$ is established, it is supposed that both $vX - wY$ and $wX + vY$ are real numbers. Therefore, it is safe to assume that $(vX-wY)^2 \geq 0$ and $(wX+vY)^2 \geq 0$, making it possible to square both sides of the latest equation from computing $Z\Omega$.
	
	
	\[ Z\Omega =  \sqrt{(vX-wY)^2 + (wX+vY)^2} = 0 \]
	\[ \Rightarrow (vX-wY)^2 + (wX+vY)^2 = 0 \]
	
	This can be reorganised by expanding and factoring out certain terms.
	\[ (vX-wY)^2 + (wX+vY)^2 = (v^2X^2 + w^2Y^2 - 2vwXY) + (w^2X^2 + v^2Y^2 + 2vwXY) \]
	\[ = (v^2 + w^2)X^2 + (v^2 + w^2)Y^2 \]
	\[ = (v^2 + w^2)(X^2+Y^2) = 0 \]
	Since $v^2 + w^2 >0$, both sides of the latest equation may be divided by $(v^2 + w^2)$.
	\[ (v^2 + w^2)(X^2+Y^2) = 0 \qquad \Rightarrow \qquad X^2+Y^2 = 0 \]
	\[ \because v, w \in \mathbb{R}, \quad v^2 + w^2 > 0 \]
	\[  \]
	
	
	
	\section{Divergence of $\zeta(\frac{2}{b})$}
	\subsection{Terms revisited}
	Now that we have established that $X^2+Y^2 = 0$ through analytic continuation, let's observe what happens. The subscripts $\zeta(s)$ will be used to denote that such terms are associated to the result of the Riemann zeta function given $s$ where $s=\frac{1}{b}+ci$.
	\[ \zeta(s) = \zeta(\frac{1}{b}+ci) := X_{\zeta(s)} + Y_{\zeta(s)}i \]
	\[ X_{\zeta(s)} = \sum\limits_{k=1}^{\infty}x_k = \sum\limits_{k=1}^{\infty}\frac{\cos({c\ln \frac{1}{k}})}{\sqrt[b]{k}}, \qquad Y_{\zeta(s)} = \sum\limits_{k=1}^{\infty}y_k = \sum\limits_{k=1}^{\infty}\frac{\sin({c\ln \frac{1}{k}})}{\sqrt[b]{k}} \]
	
	\subsection{The self-Cauchy products of X and Y}
	By definition and given all the conditions, especially $1<b<\infty$, $X_{\zeta(s)}$ and $Y_{\zeta(s)}$ diverge and therefore cannot be zero. However, $X_{\zeta(s)}^2+Y_{\zeta(s)}^2$ converges as it equals zero.
	\[ X_{\zeta(s)} \neq 0, \qquad Y_{\zeta(s)} \neq 0 \qquad\qquad (\because X_{\zeta(s)}, Y_{\zeta(s)} = \text{DNE}) \]
	\[ X_{\zeta(s)}^2+Y_{\zeta(s)}^2 = 0 \]
	
	Observe the two different Cauchy products, namely $X_{\zeta(s)}^2$ and $Y_{\zeta(s)}^2$ while we let $k \to \infty$.
	
	
	\[ \begin{array}{c}
		X_{\zeta(s)}^2 = (x_1 {+} x_2 {+} x_3 {+} \dots {+} x_k)^2 = (x_1^2 {+} x_2^2 {+} x_3^2 {+} \dots {+} x_k^2) \\
		\qquad\qquad\qquad\qquad\qquad\qquad\qquad\qquad\quad {+} 2\{(x_1x_2 {+} x_1x_3 {+} x_1x_4 {+} \dots {+} x_{1}x_{k}) \\ 
		\qquad\qquad\qquad\qquad\qquad\qquad\qquad\qquad {+} (x_2x_3 {+} x_2x_4 {+} x_2x_5 {+} \dots {+} x_2x_{k}) \\ 
		\qquad\qquad\qquad\qquad\qquad\qquad\qquad\qquad {+} (x_3x_4 {+} x_3x_5 {+} x_3x_6 {+} \dots {+} x_3x_{k}) \\ 
		\qquad\qquad\qquad\qquad\qquad\qquad\qquad\qquad \vdots \\
		\qquad\qquad\qquad\qquad {+} \dots {+} (x_{k-1}x_{k}) \} \\ 
	\end{array} \]
	\[\]
	
	\[ = \Bigl( \frac{\cos({c\ln \frac{1}{1}})}{\sqrt[b]{1}} {+} \frac{\cos({c\ln \frac{1}{2}})}{\sqrt[b]{2}} {+} \frac{\cos({c\ln \frac{1}{3}})}{\sqrt[b]{3}} {+} \dots {+} \frac{\cos({c\ln \frac{1}{k}})}{\sqrt[b]{k}} \Bigr)^2 \]
	\[ = \Bigl( \frac{1}{1} \Bigr)^{\frac2b}\cos^2({c \ln \frac{1}{1}}) {+} \Bigl( \frac{1}{2} \Bigr)^{\frac2b}\cos^2({c \ln \frac{1}{2}}) {+} \Bigl( \frac{1}{3} \Bigr)^{\frac2b}\cos^2({c \ln \frac{1}{3}}) {+} {\dots} {+} \Bigl( \frac{1}{k} \Bigr)^{\frac2b}\cos^2({c \ln \frac{1}{k}}) \]
	\[
	\begin{aligned}
		{+}2\Biggl\{ \Bigl( \frac{1}{1} \cdot \frac{1}{2}& \Bigr)^\frac1b \cos(c \ln \frac{1}{1}) \cos(c \ln \frac{1}{2}) {+} \Bigl( \frac{1}{1} \cdot \frac{1}{3} \Bigr)^\frac1b \cos(c \ln \frac{1}{1}) \cos(c \ln \frac{1}{3}) \\ 
		& {+} \Bigl( \frac{1}{1} \cdot \frac{1}{4} \Bigr)^\frac1b \cos(c \ln \frac{1}{1}) \cos(c \ln \frac{1}{4}) {+} \dots {+} \Bigl( \frac{1}{1} \cdot \frac{1}{k} \Bigr)^\frac1b \cos(c \ln \frac{1}{1}) \cos(c \ln \frac{1}{k})  \\ 		
	\end{aligned}
	\]
	\[
	\begin{aligned}
		\quad {+} \Bigl( \frac{1}{2} \cdot \frac{1}{3}& \Bigr)^\frac1b \cos(c \ln \frac{1}{2}) \cos(c \ln \frac{1}{3}) {+} \Bigl( \frac{1}{2} \cdot \frac{1}{4} \Bigr)^\frac1b \cos(c \ln \frac{1}{2}) \cos(c \ln \frac{1}{4}) \\ 
		& {+} \Bigl( \frac{1}{2} \cdot \frac{1}{5} \Bigr)^\frac1b \cos(c \ln \frac{1}{2}) \cos(c \ln \frac{1}{5}) {+} \dots {+} \Bigl( \frac{1}{2} \cdot \frac{1}{k} \Bigr)^\frac1b \cos(c \ln \frac{1}{2}) \cos(c \ln \frac{1}{k})  \\ 		
	\end{aligned}
	\]
	\[
	\begin{aligned}
		\quad {+} \Bigl( \frac{1}{3} \cdot \frac{1}{4}& \Bigr)^\frac1b \cos(c \ln \frac{1}{3}) \cos(c \ln \frac{1}{4}) {+} \Bigl( \frac{1}{3} \cdot \frac{1}{5} \Bigr)^\frac1b \cos(c \ln \frac{1}{3}) \cos(c \ln \frac{1}{5}) \\ 
		& {+} \Bigl( \frac{1}{3} \cdot \frac{1}{6} \Bigr)^\frac1b \cos(c \ln \frac{1}{3}) \cos(c \ln \frac{1}{6}) {+} \dots {+} \Bigl( \frac{1}{3} \cdot \frac{1}{k} \Bigr)^\frac1b \cos(c \ln \frac{1}{3}) \cos(c \ln \frac{1}{k}) {+} \\ 		
	\end{aligned}
	\]
	\[
	\begin{aligned}
		\vdots \qquad\qquad \vdots \qquad\qquad \vdots \qquad {+} \dots {+} \Bigl( \frac{1}{(k-1)k} \Bigr)^\frac1b \cos(c \ln \frac{1}{k-1}) \cos(c \ln \frac{1}{k}) \Biggr\}
	\end{aligned}
	\]
	\[\]
	
	\[ \begin{array}{c}
		Y_{\zeta(s)}^2 = (y_1 {+} y_2 {+} y_3 {+} \dots {+} y_k)^2 = (y_1^2 {+} y_2^2 {+} y_3^2 {+} \dots {+} y_k^2) \\
		\qquad\qquad\qquad\qquad\qquad\qquad\qquad\qquad\quad {+} 2\{(y_1y_2 {+} y_1y_3 {+} y_1y_4 {+} \dots {+} y_{1}y_{k}) \\ 
		\qquad\qquad\qquad\qquad\qquad\qquad\qquad\qquad {+} (y_2y_3 {+} y_2y_4 {+} y_2y_5 {+} \dots {+} y_2y_{k}) \\ 
		\qquad\qquad\qquad\qquad\qquad\qquad\qquad\qquad {+} (y_3y_4 {+} y_3y_5 {+} y_3y_6 {+} \dots {+} y_3y_{k}) \\ 
		\qquad\qquad\qquad\qquad\qquad\qquad\qquad\qquad \vdots \\
		\qquad\qquad\qquad\qquad {+} \dots {+} (y_{k-1}y_{k}) \} \\ 
	\end{array} \]
	\[\]
	
	\[ = \Bigl( \frac{\sin({c\ln \frac{1}{1}})}{\sqrt[b]{1}} {+} \frac{\sin({c\ln \frac{1}{2}})}{\sqrt[b]{2}} {+} \frac{\sin({c\ln \frac{1}{3}})}{\sqrt[b]{3}} {+} \dots {+} \frac{\sin({c\ln \frac{1}{k}})}{\sqrt[b]{k}} \Bigr)^2 \]
	\[ = \Bigl( \frac{1}{1} \Bigr)^{\frac2b}\sin^2({c \ln \frac{1}{1}}) {+} \Bigl( \frac{1}{2} \Bigr)^{\frac2b}\sin^2({c \ln \frac{1}{2}}) {+} \Bigl( \frac{1}{3} \Bigr)^{\frac2b}\sin^2({c \ln \frac{1}{3}}) {+} {\dots} {+} \Bigl( \frac{1}{k} \Bigr)^{\frac2b}\sin^2({c \ln \frac{1}{k}}) \]
	\[
	\begin{aligned}
		{+}2\Biggl\{ \Bigl( \frac{1}{1} \cdot \frac{1}{2}& \Bigr)^\frac1b \sin(c \ln \frac{1}{1}) \sin(c \ln \frac{1}{2}) {+} \Bigl( \frac{1}{1} \cdot \frac{1}{3} \Bigr)^\frac1b \sin(c \ln \frac{1}{1}) \sin(c \ln \frac{1}{3}) \\ 
		& {+} \Bigl( \frac{1}{1} \cdot \frac{1}{4} \Bigr)^\frac1b \sin(c \ln \frac{1}{1}) \sin(c \ln \frac{1}{4}) {+} \dots {+} \Bigl( \frac{1}{1} \cdot \frac{1}{k} \Bigr)^\frac1b \sin(c \ln \frac{1}{1}) \sin(c \ln \frac{1}{k})  \\ 		
	\end{aligned}
	\]
	\[
	\begin{aligned}
		\quad {+} \Bigl( \frac{1}{2} \cdot \frac{1}{3}& \Bigr)^\frac1b \sin(c \ln \frac{1}{2}) \sin(c \ln \frac{1}{3}) {+} \Bigl( \frac{1}{2} \cdot \frac{1}{4} \Bigr)^\frac1b \sin(c \ln \frac{1}{2}) \sin(c \ln \frac{1}{4}) \\ 
		& {+} \Bigl( \frac{1}{2} \cdot \frac{1}{5} \Bigr)^\frac1b \sin(c \ln \frac{1}{2}) \sin(c \ln \frac{1}{5}) {+} \dots {+} \Bigl( \frac{1}{2} \cdot \frac{1}{k} \Bigr)^\frac1b \sin(c \ln \frac{1}{2}) \sin(c \ln \frac{1}{k})  \\ 		
	\end{aligned}
	\]
	\[
	\begin{aligned}
		\quad {+} \Bigl( \frac{1}{3} \cdot \frac{1}{4}& \Bigr)^\frac1b \sin(c \ln \frac{1}{3}) \sin(c \ln \frac{1}{4}) {+} \Bigl( \frac{1}{3} \cdot \frac{1}{5} \Bigr)^\frac1b \sin(c \ln \frac{1}{3}) \sin(c \ln \frac{1}{5}) \\ 
		& {+} \Bigl( \frac{1}{3} \cdot \frac{1}{6} \Bigr)^\frac1b \sin(c \ln \frac{1}{3}) \sin(c \ln \frac{1}{6}) {+} \dots {+} \Bigl( \frac{1}{3} \cdot \frac{1}{k} \Bigr)^\frac1b \sin(c \ln \frac{1}{3}) \sin(c \ln \frac{1}{k}) {+} \\ 		
	\end{aligned}
	\]
	\[
	\begin{aligned}
		\vdots \qquad\qquad \vdots \qquad\qquad \vdots \qquad {+} \dots {+} \Bigl( \frac{1}{(k-1)k} \Bigr)^\frac1b \sin(c \ln \frac{1}{k-1}) \sin(c \ln \frac{1}{k}) \Biggr\}
	\end{aligned}
	\]
	\[\]
	Due to the fact that $X_{\zeta(s)}$ and $Y_{\zeta(s)}$ diverge, the two self-Cauchy products, namely $X_{\zeta(s)}^2$ and $Y_{\zeta(s)}^2$ will individually diverge. Keeping in mind that $X_{\zeta(s)}^2 + Y_{\zeta(s)}^2 = 0$ or $X_{\zeta(s)}^2 + Y_{\zeta(s)}^2$ collectively converges to zero, let the series of cosine cross terms in $X_{\zeta(s)}^2$ be $\alpha_{\zeta(s)}$ and that of sine cross terms in $Y_{\zeta(s)}^2$ be $\beta_{\zeta(s)}$.
	
	\[
	\begin{aligned}
		\alpha_{\zeta(s)} =& \, x_1x_2 + x_1x_3 + x_1x_4 + \dots + x_{1}x_{k} \\
		& + x_2x_3 + x_2x_4 + x_2x_5 + \dots + x_2x_{k} \\
		& + x_3x_4 + x_3x_5 + x_3x_6 + \dots + x_3x_{k} \\
		& \qquad\vdots \qquad\vdots \qquad\vdots \qquad\vdots \qquad\vdots \qquad\vdots \\
		&  + \dots + x_{k-1}x_{k}
	\end{aligned}
	\]
	\[\]
	\[
	\begin{aligned}
		\beta_{\zeta(s)} =& \, y_1y_2 + y_1y_3 + y_1y_4 + \dots + y_{1}y_{k} \\
		& + y_2y_3 + y_2y_4 + y_2y_5 + \dots + y_2y_{k} \\
		& + y_3y_4 + y_3y_5 + y_3y_6 + \dots + y_3y_{k} \\
		& \qquad\vdots \qquad\vdots \qquad\vdots \qquad\vdots \qquad\vdots \qquad\vdots \\
		&  + \dots + y_{k-1}y_{k}
	\end{aligned}
	\]
	\[\]
	Then we may rearrange $X_{\zeta(s)}^2$, $Y_{\zeta(s)}^2$ and $X_{\zeta(s)}^2 + Y_{\zeta(s)}^2$.
	\[ X^2_{\zeta(s)} = (x_1^2 + x_2^2 + x_3^2 + \dots + x_k^2) + 2 \alpha_{\zeta(s)} = \sum\limits_{k=1}^{\infty}x_k^2 + 2 \alpha_{\zeta(s)} \]
	
	\[ Y^2_{\zeta(s)} = (y_1^2 + y_2^2 + y_3^2 + \dots + y_k^2) + 2 \beta_{\zeta(s)} = \sum\limits_{k=1}^{\infty}y_k^2 + 2 \beta_{\zeta(s)} \]
	
	\[ X^2_{\zeta(s)} + Y^2_{\zeta(s)} = \sum\limits_{k=1}^{\infty}x_k^2 + 2\alpha_{\zeta(s)} + \sum\limits_{k=1}^{\infty}y_k^2 + 2\beta_{\zeta(s)} \]
	
	\[ = \sum\limits_{k=1}^{\infty}x_k^2 + \sum\limits_{k=1}^{\infty}y_k^2 + 2\alpha_{\zeta(s)} +  2\beta_{\zeta(s)} = 0 \]
	\[\]
	For any k, $x_k^2 + y_k^2$ can be simplified.
	\[ x_k^2 + y_k^2 = \frac{\cos^2(c \ln \frac{1}{k})}{k^{\frac{2}{b}}} + \frac{\sin^2(c \ln \frac{1}{k})}{k^{\frac{2}{b}}} \]
	\[ = \frac{\cos^2(c \ln \frac{1}{k}) + \sin^2(c \ln \frac{1}{k})}{k^{\frac{2}{b}}} = \frac{1}{k^{\frac{2}{b}}} \]
	
	
	\[\]
	As a result, the fact that $X_{\zeta(s)}^2 + Y_{\zeta(s)}^2 = 0$ leads to the formation of $\zeta(\frac{2}{b})$.
	\[ X^2_{\zeta(s)} + Y^2_{\zeta(s)} = \sum\limits_{k=1}^{\infty}\frac{1}{k^{\frac{2}{b}}} + 2(\alpha_{\zeta(s)} + \beta_{\zeta(s)}) = \zeta\Bigl( \frac{2}{b} \Bigr) + 2(\alpha_{\zeta(s)} + \beta_{\zeta(s)}) = 0 \]
	
	\[ \zeta\Bigl( \frac{2}{b} \Bigr) = -2(\alpha_{\zeta(s)} + \beta_{\zeta(s)}) \]
	
	\subsection{Divergence of $\alpha_{\zeta(s)} + \beta_{\zeta(s)}$}
	In order to study the nature of $\alpha_{\zeta(s)} + \beta_{\zeta(s)}$, it should be rearranged in a format that gives better insight.
	
	\[ \alpha_{\zeta(s)} + \beta_{\zeta(s)} = \qquad\qquad\qquad\qquad\qquad\qquad\qquad\qquad\qquad\qquad\qquad\qquad (k \to \infty) \]
	\[
	\begin{aligned}
		& \Bigl( \frac{1}{1} \cdot \frac{1}{2} \Bigr)^\frac1b \cos(c \ln \frac{1}{1}) \cos(c \ln \frac{1}{2}) + \Bigl( \frac{1}{1} \cdot \frac{1}{3} \Bigr)^\frac1b \cos(c \ln \frac{1}{1}) \cos(c \ln \frac{1}{3}) + \\
		& \Bigl( \frac{1}{1} \cdot \frac{1}{2} \Bigr)^\frac1b \sin(c \ln \frac{1}{1}) \sin(c \ln \frac{1}{2}) + \Bigl( \frac{1}{1} \cdot \frac{1}{3} \Bigr)^\frac1b \sin(c \ln \frac{1}{1}) \sin(c \ln \frac{1}{3}) + 
	\end{aligned}
	\]
	
	\[
	\begin{aligned}
		\qquad\qquad\quad & \Bigl( \frac{1}{1} \cdot \frac{1}{4} \Bigr)^\frac1b \cos(c \ln \frac{1}{1}) \cos(c \ln \frac{1}{4}) + \dots + \Bigl( \frac{1}{1} \cdot \frac{1}{k} \Bigr)^\frac1b \cos(c \ln \frac{1}{1}) \cos(c \ln \frac{1}{k}) + \\
		& \Bigl( \frac{1}{1} \cdot \frac{1}{4} \Bigr)^\frac1b \sin(c \ln \frac{1}{1}) \sin(c \ln \frac{1}{4}) + \dots + \Bigl( \frac{1}{1} \cdot \frac{1}{k} \Bigr)^\frac1b \sin(c \ln \frac{1}{1}) \sin(c \ln \frac{1}{k}) +
	\end{aligned}
	\]
	\[
	\begin{aligned}
		& \Bigl( \frac{1}{2} \cdot \frac{1}{3} \Bigr)^\frac1b \cos(c \ln \frac{1}{2}) \cos(c \ln \frac{1}{3}) + \Bigl( \frac{1}{2} \cdot \frac{1}{4} \Bigr)^\frac1b \cos(c \ln \frac{1}{2}) \cos(c \ln \frac{1}{4}) + \\
		& \Bigl( \frac{1}{2} \cdot \frac{1}{3} \Bigr)^\frac1b \sin(c \ln \frac{1}{2}) \sin(c \ln \frac{1}{3}) + \Bigl( \frac{1}{2} \cdot \frac{1}{4} \Bigr)^\frac1b \sin(c \ln \frac{1}{2}) \sin(c \ln \frac{1}{4}) + 
	\end{aligned}
	\]
	\[
	\begin{aligned}
		\qquad\qquad\quad & \Bigl( \frac{1}{2} \cdot \frac{1}{5} \Bigr)^\frac1b \cos(c \ln \frac{1}{2}) \cos(c \ln \frac{1}{5}) + \dots + \Bigl( \frac{1}{2} \cdot \frac{1}{k} \Bigr)^\frac1b \cos(c \ln \frac{1}{2}) \cos(c \ln \frac{1}{k}) + \\
		& \Bigl( \frac{1}{2} \cdot \frac{1}{5} \Bigr)^\frac1b \sin(c \ln \frac{1}{2}) \sin(c \ln \frac{1}{5}) + \dots + \Bigl( \frac{1}{2} \cdot \frac{1}{k} \Bigr)^\frac1b \sin(c \ln \frac{1}{2}) \sin(c \ln \frac{1}{k}) +
	\end{aligned}
	\]
	\[
	\begin{aligned}
		& \Bigl( \frac{1}{3} \cdot \frac{1}{4} \Bigr)^\frac1b \cos(c \ln \frac{1}{3}) \cos(c \ln\frac{1}{4}) + \Bigl( \frac{1}{3} \cdot \frac{1}{5} \Bigr)^\frac1b \cos(c \ln \frac{1}{3}) \cos(c \ln \frac{1}{5}) + \\
		& \Bigl( \frac{1}{3} \cdot \frac{1}{4} \Bigr)^\frac1b \sin(c \ln \frac{1}{3}) \sin(c \ln \frac{1}{4}) + \Bigl( \frac{1}{3} \cdot \frac{1}{5} \Bigr)^\frac1b \sin(c \ln \frac{1}{3}) \sin(c \ln \frac{1}{5}) + 
	\end{aligned}
	\]
	\[
	\begin{aligned}
		\qquad\qquad\quad & \Bigl( \frac{1}{3} \cdot \frac{1}{6} \Bigr)^\frac1b \cos(c \ln \frac{1}{3}) \cos(c \ln \frac{1}{6}) + \dots + \Bigl( \frac{1}{3} \cdot \frac{1}{k} \Bigr)^\frac1b \cos(c \ln \frac{1}{3}) \cos(c \ln \frac{1}{k}) + \\
		& \Bigl( \frac{1}{3} \cdot \frac{1}{6} \Bigr)^\frac1b \sin(c \ln \frac{1}{3}) \sin(c \ln \frac{1}{6}) + \dots + \Bigl( \frac{1}{3} \cdot \frac{1}{k} \Bigr)^\frac1b \sin(c \ln \frac{1}{3}) \sin(c \ln \frac{1}{k}) +
	\end{aligned}
	\]
	
	\[ \qquad\qquad \dots \qquad + \Bigl( \frac{1}{(k-1)k} \Bigr)^\frac1b \cos(c \ln \frac{1}{k-1}) \cos(c \ln \frac{1}{k}) \]
	\[ \qquad\qquad \dots \qquad + \Bigl( \frac{1}{(k-1)k} \Bigr)^\frac1b \sin(c \ln \frac{1}{k-1}) \sin(c \ln \frac{1}{k}) \]
	
	\[\]
	By initial grouping, we can have similar cosine and sine terms together.
	\[
	\begin{aligned}
		= \quad & \Bigl( \frac{1}{1} \cdot \frac{1}{2} \Bigr)^\frac1b \Bigl( \cos(c \ln \frac{1}{1}) \cos(c \ln \frac{1}{2}) + \sin(c \ln \frac{1}{1}) \sin(c \ln \frac{1}{2}) \Bigr) + \\
		& \qquad \Bigl( \frac{1}{1} \cdot \frac{1}{3} \Bigr)^\frac1b \Bigl( \cos(c \ln \frac{1}{1}) \cos(c \ln \frac{1}{3}) + \sin(c \ln \frac{1}{1}) \sin(c \ln \frac{1}{3}) \Bigr) + \\
		& \qquad \Bigl( \frac{1}{1} \cdot \frac{1}{4} \Bigr)^\frac1b \Bigl( \cos(c \ln \frac{1}{1}) \cos(c \ln \frac{1}{4}) + \sin(c \ln \frac{1}{1}) \sin(c \ln \frac{1}{4}) \Bigr) + \\
		& \qquad \dots  + \Bigl( \frac{1}{1} \cdot \frac{1}{k} \Bigr)^\frac1b \Bigl( \cos(c \ln \frac{1}{1}) \cos(c \ln \frac{1}{k}) + \sin(c \ln \frac{1}{1}) \sin(c \ln \frac{1}{k}) \Bigr) + 
	\end{aligned}
	\]
	\[
	\begin{aligned}
		\quad & \Bigl( \frac{1}{2} \cdot \frac{1}{3} \Bigr)^\frac1b \Bigl( \cos(c \ln \frac{1}{2}) \cos(c \ln \frac{1}{3}) + \sin(c \ln \frac{1}{2}) \sin(c \ln \frac{1}{3}) \Bigr) + \\
		& \qquad \Bigl( \frac{1}{2} \cdot \frac{1}{4} \Bigr)^\frac1b \Bigl( \cos(c \ln \frac{1}{2}) \cos(c \ln \frac{1}{4}) + \sin(c \ln \frac{1}{2}) \sin(c \ln \frac{1}{4}) \Bigr) + \\
		& \qquad \Bigl( \frac{1}{2} \cdot \frac{1}{5} \Bigr)^\frac1b \Bigl( \cos(c \ln \frac{1}{2}) \cos(c \ln \frac{1}{5}) + \sin(c \ln \frac{1}{2}) \sin(c \ln \frac{1}{5}) \Bigr) + \\
		& \qquad \dots  + \Bigl( \frac{1}{2} \cdot \frac{1}{k} \Bigr)^\frac1b \Bigl( \cos(c \ln \frac{1}{2}) \cos(c \ln \frac{1}{k}) + \sin(c \ln \frac{1}{2}) \sin(c \ln \frac{1}{k}) \Bigr) + 
	\end{aligned}
	\]
	\[
	\begin{aligned}
		\quad & \Bigl( \frac{1}{3} \cdot \frac{1}{4} \Bigr)^\frac1b \Bigl( \cos(c \ln \frac{1}{3}) \cos(c \ln \frac{1}{3}) + \sin(c \ln \frac{1}{3}) \sin(c \ln \frac{1}{3}) \Bigr) + \\
		& \qquad \Bigl( \frac{1}{3} \cdot \frac{1}{5} \Bigr)^\frac1b \Bigl( \cos(c \ln \frac{1}{3}) \cos(c \ln \frac{1}{5}) + \sin(c \ln \frac{1}{3}) \sin(c \ln \frac{1}{5}) \Bigr) + \\
		& \qquad \Bigl( \frac{1}{3} \cdot \frac{1}{6} \Bigr)^\frac1b \Bigl( \cos(c \ln \frac{1}{3}) \cos(c \ln \frac{1}{6}) + \sin(c \ln \frac{1}{3}) \sin(c \ln \frac{1}{6}) \Bigr) + \\
		& \qquad \dots  + \Bigl( \frac{1}{3} \cdot \frac{1}{k} \Bigr)^\frac1b \Bigl( \cos(c \ln \frac{1}{3}) \cos(c \ln \frac{1}{k}) + \sin(c \ln \frac{1}{3}) \sin(c \ln \frac{1}{k}) \Bigr) + 
	\end{aligned}
	\]
	\[ \qquad\qquad \dots + \Bigl( \frac{1}{(k-1)k} \Bigr)^\frac1b \Bigl( \cos(c \ln \frac{1}{k-1}) \cos(c \ln \frac{1}{k}) + \sin(c \ln \frac{1}{k-1}) \sin(c \ln \frac{1}{k}) \Bigr) \]
	
	\[\]
	Then we may utilise the fact that $ \cos{A} \cdot \cos{B} = \cos{(A-B)} - \sin{A}\sin{B} $ to cancel out all the sine terms.
	
	\[
	\begin{aligned}
		= \quad & \Bigl( \frac{1}{1} \cdot \frac{1}{2} \Bigr)^\frac1b \Bigl( \cos(c \ln \frac{1}{1} - c \ln \frac{1}{2}) - \sin(c \ln \frac{1}{1}) \sin(c \ln \frac{1}{2}) + \sin(c \ln \frac{1}{1}) \sin(c \ln \frac{1}{2}) \Bigr) + \\
		& \qquad \Bigl( \frac{1}{1} \cdot \frac{1}{3} \Bigr)^\frac1b \Bigl( \cos(c \ln \frac{1}{1} - c \ln \frac{1}{3}) - \sin(c \ln \frac{1}{1}) \sin(c \ln \frac{1}{3}) + \sin(c \ln \frac{1}{1}) \sin(c \ln \frac{1}{3}) \Bigr) + \\
		& \qquad \Bigl( \frac{1}{1} \cdot \frac{1}{4} \Bigr)^\frac1b \Bigl( \cos(c \ln \frac{1}{1} - c \ln \frac{1}{4}) - \sin(c \ln \frac{1}{1}) \sin(c \ln \frac{1}{4}) + \sin(c \ln \frac{1}{1}) \sin(c \ln \frac{1}{4}) \Bigr) + \\
		& \dots  + \Bigl( \frac{1}{1} \cdot \frac{1}{k} \Bigr)^\frac1b \Bigl( \cos(c \ln \frac{1}{1} - c \ln \frac{1}{k}) - \sin(c \ln \frac{1}{1}) \sin(c \ln \frac{1}{k}) + \sin(c \ln \frac{1}{1}) \sin(c \ln \frac{1}{k}) \Bigr) + 
	\end{aligned}
	\]
	\[
	\begin{aligned}
		\quad & \Bigl( \frac{1}{2} \cdot \frac{1}{3} \Bigr)^\frac1b \Bigl( \cos(c \ln \frac{1}{2} - c \ln \frac{1}{3}) - \sin(c \ln \frac{1}{2}) \sin(c \ln \frac{1}{3}) + \sin(c \ln \frac{1}{2}) \sin(c \ln \frac{1}{3}) \Bigr) + \\
		& \qquad \Bigl( \frac{1}{2} \cdot \frac{1}{4} \Bigr)^\frac1b \Bigl( \cos(c \ln \frac{1}{2} - c \ln \frac{1}{4}) - \sin(c \ln \frac{1}{2}) \sin(c \ln \frac{1}{4}) + \sin(c \ln \frac{1}{2}) \sin(c \ln \frac{1}{4}) \Bigr) + \\
		& \qquad \Bigl( \frac{1}{2} \cdot \frac{1}{5} \Bigr)^\frac1b \Bigl( \cos(c \ln \frac{1}{2} - c \ln \frac{1}{5}) - \sin(c \ln \frac{1}{2}) \sin(c \ln \frac{1}{5}) + \sin(c \ln \frac{1}{2}) \sin(c \ln \frac{1}{5}) \Bigr) + \\
		& \dots  + \Bigl( \frac{1}{2} \cdot \frac{1}{k} \Bigr)^\frac1b \Bigl( \cos(c \ln \frac{1}{2} - c \ln \frac{1}{k}) - \sin(c \ln \frac{1}{2}) \sin(c \ln \frac{1}{k}) + \sin(c \ln \frac{1}{2}) \sin(c \ln \frac{1}{k}) \Bigr) + 
	\end{aligned}
	\]
	\[
	\begin{aligned}
		\quad & \Bigl( \frac{1}{3} \cdot \frac{1}{4} \Bigr)^\frac1b \Bigl( \cos(c \ln \frac{1}{3} - c \ln \frac{1}{4}) - \sin(c \ln \frac{1}{3}) \sin(c \ln \frac{1}{4}) + \sin(c \ln \frac{1}{3}) \sin(c \ln \frac{1}{4}) \Bigr) + \\
		& \qquad \Bigl( \frac{1}{3} \cdot \frac{1}{5} \Bigr)^\frac1b \Bigl( \cos(c \ln \frac{1}{3} - c \ln \frac{1}{5}) - \sin(c \ln \frac{1}{3}) \sin(c \ln \frac{1}{5}) + \sin(c \ln \frac{1}{3}) \sin(c \ln \frac{1}{5}) \Bigr) + \\
		& \qquad \Bigl( \frac{1}{3} \cdot \frac{1}{6} \Bigr)^\frac1b \Bigl( \cos(c \ln \frac{1}{3} - c \ln \frac{1}{6}) - \sin(c \ln \frac{1}{3}) \sin(c \ln \frac{1}{6}) + \sin(c \ln \frac{1}{3}) \sin(c \ln \frac{1}{6}) \Bigr) + \\
		& \dots  + \Bigl( \frac{1}{3} \cdot \frac{1}{k} \Bigr)^\frac1b \Bigl( \cos(c \ln \frac{1}{3} - c \ln \frac{1}{k}) - \sin(c \ln \frac{1}{3}) \sin(c \ln \frac{1}{k}) + \sin(c \ln \frac{1}{3}) \sin(c \ln \frac{1}{k}) \Bigr) + 
	\end{aligned}
	\]
	%\[ \dots + \Bigl( \frac{1}{(k-1)k} \Bigr)^\frac1b \Bigl( \cos(c \ln \frac{1}{k-1} - c \ln \frac{1}{k}) - \sin(c \ln \frac{1}{k-1}) \sin(c \ln \frac{1}{k}) + \sin(c \ln \frac{1}{k-1}) \sin(c \ln \frac{1}{k}) \Bigr) \]
	$ \dots + \Bigl( \frac{1}{(k-1)k} \Bigr)^\frac1b \Bigl( \cos(c \ln \frac{1}{k-1} - c \ln \frac{1}{k}) - \sin(c \ln \frac{1}{k-1}) \sin(c \ln \frac{1}{k}) + \sin(c \ln \frac{1}{k-1}) \sin(c \ln \frac{1}{k}) \Bigr) $
	
	\[
	\begin{aligned}
		= \quad & \Bigl( \frac{1}{1} \cdot \frac{1}{2} \Bigr)^\frac1b \Bigl( \cos(c \ln \frac{1}{1} - c \ln \frac{1}{2}) \Bigr) + \Bigl( \frac{1}{1} \cdot \frac{1}{3} \Bigr)^\frac1b \Bigl( \cos(c \ln \frac{1}{1} - c \ln \frac{1}{3}) \Bigr) + \\
		& \qquad \Bigl( \frac{1}{1} \cdot \frac{1}{4} \Bigr)^\frac1b \Bigl( \cos(c \ln \frac{1}{1} - c \ln \frac{1}{4}) \Bigr) + \dots  + \Bigl( \frac{1}{1} \cdot \frac{1}{k} \Bigr)^\frac1b \Bigl( \cos(c \ln \frac{1}{1} - c \ln \frac{1}{k}) \Bigr) + 
	\end{aligned}
	\]
	\[
	\begin{aligned}
		\quad\quad & \Bigl( \frac{1}{2} \cdot \frac{1}{3} \Bigr)^\frac1b \Bigl( \cos(c \ln \frac{1}{2} - c \ln \frac{1}{3}) \Bigr) + \Bigl( \frac{1}{2} \cdot \frac{1}{4} \Bigr)^\frac1b \Bigl( \cos(c \ln \frac{1}{2} - c \ln \frac{1}{4}) \Bigr) + \\
		& \qquad \Bigl( \frac{1}{2} \cdot \frac{1}{5} \Bigr)^\frac1b \Bigl( \cos(c \ln \frac{1}{2} - c \ln\frac{1}{5}) \Bigr) + \dots  + \Bigl( \frac{1}{2} \cdot \frac{1}{k} \Bigr)^\frac1b \Bigl( \cos(c \ln \frac{1}{2} - c \ln \frac{1}{k}) \Bigr) + 
	\end{aligned}
	\]
	\[
	\begin{aligned}
		\quad\quad & \Bigl( \frac{1}{3} \cdot \frac{1}{4} \Bigr)^\frac1b \Bigl( \cos(c \ln \frac{1}{3} - c \ln \frac{1}{4}) \Bigr) + \Bigl( \frac{1}{3} \cdot \frac{1}{5} \Bigr)^\frac1b \Bigl( \cos(c \ln \frac{1}{3} - c \ln \frac{1}{5}) \Bigr) + \\
		& \qquad \Bigl( \frac{1}{3} \cdot \frac{1}{6} \Bigr)^\frac1b \Bigl( \cos(c \ln \frac{1}{3} - c \ln \frac{1}{6}) \Bigr) + \dots  + \Bigl( \frac{1}{3} \cdot \frac{1}{k} \Bigr)^\frac1b \Bigl( \cos(c \ln \frac{1}{3} - c \ln \frac{1}{k}) \Bigr) + 
	\end{aligned}
	\]
	\[ \dots + \Bigl( \frac{1}{(k-1)k} \Bigr)^\frac1b \Bigl( \cos(c \ln \frac{1}{k-1} - c \ln \frac{1}{k}) \Bigr) \]
	\[\]
	
	Common terms of each group can be factored out.
	\[
	\begin{aligned}
		= \quad & \Bigl( \frac{1}{1} \Bigr)^\frac1b \Biggl\{ \Bigl( \frac{1}{2} \Bigr)^\frac1b \Bigl( \cos(c \ln \frac{1}{1} - c \ln \frac{1}{2}) \Bigr) + \Bigl( \frac{1}{3} \Bigr)^\frac1b \Bigl( \cos(c \ln \frac{1}{1} - c \ln \frac{1}{3}) \Bigr) + \\
		& \qquad \Bigl( \frac{1}{4} \Bigr)^\frac1b \Bigl( \cos(c \ln \frac{1}{1} - c \ln \frac{1}{4}) \Bigr) + \dots  + \Bigl( \frac{1}{k} \Bigr)^\frac1b \Bigl( \cos(c \ln \frac{1}{1} - c \ln \frac{1}{k}) \Bigr) \Biggr\} + 
	\end{aligned}
	\]
	\[
	\begin{aligned}
		\quad\quad & \Bigl( \frac{1}{2}  \Bigr)^\frac1b \Biggl\{ \Bigl( \frac{1}{3} \Bigr)^\frac1b \Bigl( \cos(c \ln \frac{1}{2} - c \ln \frac{1}{3}) \Bigr) + \Bigl( \frac{1}{4} \Bigr)^\frac1b \Bigl( \cos(c \ln \frac{1}{2} - c \ln \frac{1}{4}) \Bigr) + \\
		& \qquad \Bigl( \frac{1}{5} \Bigr)^\frac1b \Bigl( \cos(c \ln \frac{1}{2} - c \ln \frac{1}{5}) \Bigr) + \dots  + \Bigl( \frac{1}{k} \Bigr)^\frac1b \Bigl( \cos(c \ln \frac{1}{2} - c \ln \frac{1}{k}) \Bigr) \Biggr\} + 
	\end{aligned}
	\]
	\[
	\begin{aligned}
		\quad\quad & \Bigl( \frac{1}{3} \Bigr)^\frac1b \Biggl\{ \Bigl( \frac{1}{4} \Bigr)^\frac1b \Bigl( \cos(c \ln \frac{1}{3} - c \ln \frac{1}{4}) \Bigr) + \Bigl( \frac{1}{5} \Bigr)^\frac1b \Bigl( \cos(c \ln \frac{1}{3} - c \ln \frac{1}{5}) \Bigr) + \\
		& \qquad \Bigl( \frac{1}{6} \Bigr)^\frac1b \Bigl( \cos(c \ln \frac{1}{3} - c \ln \frac{1}{6}) \Bigr) + \dots  + \Bigl( \frac{1}{k} \Bigr)^\frac1b \Bigl( \cos(c \ln \frac{1}{3} - c \ln \frac{1}{k}) \Bigr) \Biggr\} + 
	\end{aligned}
	\]
	\[ \dots + \Bigl( \frac{1}{k-1} \Bigr)^\frac1b \Biggl\{ \Bigl( \frac{1}{k} \Bigr)^\frac1b \cos(c \ln \frac{1}{k-1} - c \ln \frac{1}{k}) \Biggr\} \]
	
	
	Finally, the logarithm trick where $ \ln{A} - \ln{B} = \ln{\frac{A}{B}} $ will merge all the $c \cdot \ln \frac{1}{k}$ wherever possible.
	
	\[
	\begin{aligned}
		= \quad \Bigl( \frac{1}{1} \Bigr)^\frac1b \Biggl\{ \Bigl( \frac{1}{2} \Bigr)^\frac1b & \Bigl( \cos(c \ln \frac{2}{1}) \Bigr) + \Bigl( \frac{1}{3} \Bigr)^\frac1b \Bigl( \cos(c \ln \frac{3}{1}) \Bigr) + \\
		& \Bigl( \frac{1}{4} \Bigr)^\frac1b \Bigl( \cos(c \ln \frac{4}{1}) \Bigr) + \dots  + \Bigl( \frac{1}{k} \Bigr)^\frac1b \Bigl( \cos(c \ln \frac{k}{1}) \Bigr) \Biggr\} + 
	\end{aligned}
	\]
	\[
	\begin{aligned}
		\quad\quad \Bigl( \frac{1}{2}  \Bigr)^\frac1b \Biggl\{ \Bigl( \frac{1}{3} \Bigr)^\frac1b & \Bigl( \cos(c \ln \frac{3}{2}) \Bigr) + \Bigl( \frac{1}{4} \Bigr)^\frac1b \Bigl( \cos(c \ln \frac{4}{2}) \Bigr) + \\
		& \Bigl( \frac{1}{5} \Bigr)^\frac1b \Bigl( \cos(c \ln \frac{5}{2}) \Bigr) + \dots  + \Bigl( \frac{1}{k} \Bigr)^\frac1b \Bigl( \cos(c \ln \frac{k}{2}) \Bigr) \Biggr\} + 
	\end{aligned}
	\]
	\[
	\begin{aligned}
		\quad\quad \Bigl( \frac{1}{3} \Bigr)^\frac1b \Biggl\{ \Bigl( \frac{1}{4} \Bigr)^\frac1b & \Bigl( \cos(c \ln \frac{4}{3}) \Bigr) + \Bigl( \frac{1}{5} \Bigr)^\frac1b \Bigl( \cos(c \ln \frac{5}{3}) \Bigr) + \\
		& \Bigl( \frac{1}{6} \Bigr)^\frac1b \Bigl( \cos(c \ln \frac{6}{3}) \Bigr) + \dots  + \Bigl( \frac{1}{k} \Bigr)^\frac1b \Bigl( \cos(c \ln \frac{k}{3}) \Bigr) \Biggr\} + 
	\end{aligned}
	\]
	\[ \dots + \Bigl( \frac{1}{k-1} \Bigr)^\frac1b \Biggl\{ \Bigl( \frac{1}{k} \Bigr)^\frac1b \cos(c \ln \frac{k}{k-1}) \Biggr\} \]
	
	\[
	\begin{aligned}
		= \quad \Bigl( \frac{1}{1} \Bigr)^\frac1b \Biggl\{ \Bigl( \frac{1}{2} \Bigr)^\frac1b & \Bigl( \cos(c \ln \frac{1}{2}) \Bigr) + \Bigl( \frac{1}{3} \Bigr)^\frac1b \Bigl( \cos(c \ln \frac{1}{3}) \Bigr) + \\
		& \Bigl( \frac{1}{4} \Bigr)^\frac1b \Bigl( \cos(c \ln \frac{1}{4}) \Bigr) + \dots  + \Bigl( \frac{1}{k} \Bigr)^\frac1b \Bigl( \cos(c \ln \frac{1}{k}) \Bigr) \Biggr\} + 
	\end{aligned}
	\]
	\[
	\begin{aligned}
		\quad\quad \Bigl( \frac{1}{2}  \Bigr)^\frac1b \Biggl\{ \Bigl( \frac{1}{3} \Bigr)^\frac1b & \Bigl( \cos(c \ln \frac{2}{3}) \Bigr) + \Bigl( \frac{1}{4} \Bigr)^\frac1b \Bigl( \cos(c \ln \frac{2}{4}) \Bigr) + \\
		& \Bigl( \frac{1}{5} \Bigr)^\frac1b \Bigl( \cos(c \ln \frac{2}{5}) \Bigr) + \dots  + \Bigl( \frac{1}{k} \Bigr)^\frac1b \Bigl( \cos(c \ln \frac{2}{k}) \Bigr) \Biggr\} + 
	\end{aligned}
	\]
	\[
	\begin{aligned}
		\quad\quad \Bigl( \frac{1}{3} \Bigr)^\frac1b \Biggl\{ \Bigl( \frac{1}{4} \Bigr)^\frac1b & \Bigl( \cos(c \ln \frac{3}{4}) \Bigr) + \Bigl( \frac{1}{5} \Bigr)^\frac1b \Bigl( \cos(c \ln \frac{3}{5}) \Bigr) + \\
		& \Bigl( \frac{1}{6} \Bigr)^\frac1b \Bigl( \cos(c \ln \frac{3}{6}) \Bigr) + \dots  + \Bigl( \frac{1}{k} \Bigr)^\frac1b \Bigl( \cos(c \ln \frac{3}{k}) \Bigr) \Biggr\} + 
	\end{aligned}
	\]
	\[ \dots + \Bigl( \frac{1}{k-1} \Bigr)^\frac1b \Biggl\{ \Bigl( \frac{1}{k} \Bigr)^\frac1b \cos(c \ln \frac{k-1}{k}) \Biggr\} \]
	
	
	
	
	\[\]
	
	The reorganised $\alpha_{\zeta(s)} + \beta_{\zeta(s)}$ can be viewed as a collection of nodes and branches where each node is the outer coefficient in the form of $\frac{1}{m}^{\frac{1}{b}}$ where $m \in \mathbb{N}$ and each branch is a series of cosine terms with their own inner coefficients. Neither the nodes nor the branches converge. The nodes form a series that would be defined as $\zeta(\frac{1}{b})$ and given that $0 < \frac{1}{b} < 1$, this series diverges. The branches, where each one is a sequence of cosine values multiplied by summands from a subseries of $\zeta(\frac{1}{b})$ diverge. The cosine values as numerators oscillate between -1 and 1, at an increasingly slower rate as $k$ per branch tends to infinity, including both bounds. The terms from $\zeta(\frac{1}{b})$ within every branch act as denominators when multiplied and do not diminish at a sufficient rate to cause each branch to converge. As a result, each node-branch product will diverge, rendering the entire series, a nested loop of Dirichlet series, into divergence. If $\alpha_{\zeta(s)} + \beta_{\zeta(s)}$ diverges, $l {\cdot} (\alpha_{\zeta(s)} + \beta_{\zeta(s)})$ diverges where $l \in \mathbb{R}_{\neq 0}$. This means that $\zeta(\frac{2}{b})$ diverges as $\zeta(\frac{2}{b}) = -2(\alpha_{\zeta(s)} + \beta_{\zeta(s)})$. 
	
	\[ \zeta\Bigl( \frac{2}{b} \Bigr) \quad \text{DNE} \qquad (0 < \frac{1}{b} < 1) \]
	
	
	\section{Analytically continued $\zeta(\frac{2}{b})$}
	\subsection{Additional analytical continuation}
	Although $\zeta(\frac{2}{b})$ and $\alpha_{\zeta(s)} + \beta_{\zeta(s)}$ diverge, it is premature to assume that $b=2$ in an attempt to claim that $\zeta(\frac{2}{b})$ is the single pole, $\zeta(1)$. Since the Riemann zeta function, without analytic continuation, diverges when $Re(s) \leq 1$, the divergent property of the equation $\zeta(\frac{2}{b}) = -2(\alpha_{\zeta(s)} + \beta_{\zeta(s)})$ may simply indicate that $\frac{2}{b} \leq 1$. This is within the initially defined range of $\frac{1}{b}$ as $0 < \frac{1}{b} < 1$ or $0 < \frac{2}{b} < 2$.
	
	Analytic continuation of $\zeta(\frac{2}{b})$, eventually expressed as $\frac{\eta(\frac{2}{b})}{D_{2/b}}$, may help us understand more about $\zeta(\frac{2}{b})$, $\frac{\eta(\frac{2}{b})}{D_{2/b}}$ and another variant of $\alpha + \beta$. However, simply dividing both sides by $D_{2/b} = 1 - 2^{1-\frac{2}{b}}$ does not produce the appropriate results as $\alpha_{\zeta(s)} + \beta_{\zeta(s)}$ originates from a naïve form of Riemann zeta function given $s$. In order to obtain its analytically continued version, $\alpha_{\eta(s)} + \beta_{\eta(s)}$, a similar set of processes adopting the same concepts of Cartesian coordinate system and Euclidean distance have to be executed.
	
	\subsection{Expressing $\eta(s)$ and its summands}
	Instead of $\zeta(s)$, $\eta(s)$ will be the new main focus. In similar fashion, each term of the series from $\eta(s)$ can be expressed in $b$ and $c$.
	\[ \frac{(-1)^{k-1}}{k^s} = \frac{(-1)^{k-1}}{k^{\frac{1}{b} + ci}} = \frac{\cos({c\ln k})(-1)^{k-1}}{\sqrt[b]{k}} + \frac{-\sin({c\ln k})(-1)^{k-1}}{\sqrt[b]{k}} i \]
	Then, a new $z$ can be set to equal $\eta(s) = G + Hi$. Here, $z$ may again be a point on an XY plane.
	\[ \eta(s) \qquad = \qquad +\frac{1}{1^s} \qquad \Big|\Big|  \qquad +\frac{\cos({c\ln 1})}{\sqrt[b]{1}} \qquad +\frac{-\sin({c\ln 1})}{\sqrt[b]{1}} \quad i  \]
	
	\[ \quad \qquad  \qquad \qquad -\frac{1}{2^s} \qquad \Big|\Big|  \qquad -\frac{\cos({c\ln 2})}{\sqrt[b]{2}} \qquad -\frac{-\sin({c\ln 2})}{\sqrt[b]{2}} \quad i  \]
	
	\[ \quad \qquad  \qquad \qquad +\frac{1}{3^s} \qquad \Big|\Big|  \qquad +\frac{\cos({c\ln 3})}{\sqrt[b]{3}} \qquad +\frac{-\sin({c\ln 3})}{\sqrt[b]{3}} \quad i  \]
	
	\[ \vdots \]
	
	\[ \qquad  \quad +\frac{(-1)^{k-1}}{k^s} \qquad \Big|\Big|  \qquad \frac{\cos({c\ln k}){\cdot(-1)^{k-1}}}{\sqrt[b]{k}} \qquad \frac{{-}\sin({c\ln k}){\cdot({-}1)^{k-1}}}{\sqrt[b]{k}} \quad i  \]
	
	\[\quad \rule{12 cm}{.8pt} \quad \]
	\[ \qquad \qquad \qquad z \qquad\qquad =  \qquad\qquad G \qquad + \qquad H \quad i \qquad\qquad \]
	
	\[ \text{or} \]
	\[ z (G, H) \]
	\[\]
	
	If $G$ and $H$ are series of $g_k$ and $h_k$, respectively, we may reuse $x_k$ and $y_k$ to express these new terms.
	\[ G = Re(\eta(s)) = \sum\limits_{k=1}^{\infty}g^k, \quad g_k = (-1)^{k-1} \cdot x_k \]
	\[ H = Im(\eta(s)) = \sum\limits_{k=1}^{\infty}h^k, \quad h_k = (-1)^{k-1} \cdot y_k \]
	
	Now $\zeta(s)$ may be expressed in terms of $G$ and $H$ to obtain the a $Z$.
	\[ \zeta(s) = \frac{\eta(s)}{D_{s}} = \frac{G + Hi}{1-2^{1-s}} = \frac{(G+Hi)(v-wi)}{(v+wi)(v-wi)} \]
	
	\[ = \frac{vG-wGi+vHi-wHi^2}{v^2+w^2} = \frac{(vG+wH)+(-wG+vH)i}{v^2+w^2} \]
	
	\[ = \frac{vG+wH}{v^2+w^2} + \frac{(-wG+vH)}{v^2+w^2}i \]
	
	\[ \text{or} \]
	
	\[ \quad Z\biggr(\frac{vG+wH}{v^2+w^2}, \quad \frac{-wG+vH}{v^2+w^2}\biggl) \]
	
	\subsection{The new $Z\Omega$}
	The distance between the same arbitrary point $\Omega(p, q)$ and the new $Z$ will be calculated to further observe properties of $G$ and $H$.
	\[ Z\Omega = \sqrt{\Bigl( \frac{vG+wH}{v^2+w^2}-p \Bigr)^2 + \Bigl( \frac{-wG+vH}{v^2+w^2}-q \Bigr)^2} \]
	\[\]
	$ \text{If} \quad Z\Omega = 0, $
	\[ Z\Omega = \sqrt{\Bigl( \frac{vG+wH}{v^2+w^2}-p \Bigr)^2 + \Bigl( \frac{-wG+vH}{v^2+w^2}-q \Bigr)^2} = 0 \]
	
	\[ = \sqrt{  (\frac{vG{+}wH}{v^2{+}w^2})^2 {+} p^2 {-} 2p(\frac{vG{+}wH}{v^2{+}w^2}) {+} (\frac{-wG{+}vH}{v^2{+}w^2})^2 {+} q^2 {-}2q(\frac{-wG{+}vH}{v^2{+}w^2}) } = 0 \]
	\[\]
	\[\]
	$ \text{If} \quad \Omega(p, q) = (0, 0), $
	\[ Z\Omega = \sqrt{  (\frac{vG{+}wH}{v^2{+}w^2})^2 {+} 0^2 {-} 2 {\cdot} 0 (\frac{vG{+}wH}{v^2{+}w^2}) {+} (\frac{-wG{+}vH}{v^2{+}w^2})^2 {+} 0^2 {-} 2 {\cdot} 0 (\frac{-wG{+}vH}{v^2{+}w^2}) } \]
	
	\[ = \sqrt{\Bigl(\frac{vG{+}wH}{v^2{+}w^2}\Bigr)^2 + \Bigl( \frac{-wG{+}vH}{v^2{+}w^2} \Bigr)^2} = 0 \]
	
	We square both sides and expand for $G^2$ and $H^2$.
	
	\[ \Bigl(\frac{vG{+}wH}{v^2{+}w^2}\Bigr)^2 + \Bigl(\frac{-wG{+}vH}{v^2{+}w^2}\Bigr)^2 = 0 \]
	
	\[ = \frac{v^2G^2 + w^2H^2 + 2vwGH}{(v^2+w^2)^2} + \frac{w^2G^2 + v^2H^2 - 2vwGH}{(v^2+w^2)^2} \]
	\[ = \frac{v^2G^2 + w^2H^2 + w^2G^2 + v^2H^2}{(v^2+w^2)^2} \]
	\[ = \frac{(v^2+w^2)G^2 + (v^2+w^2)H^2}{(v^2+w^2)^2} \]
	\[ = \frac{(v^2+w^2)(G^2 + H^2)}{(v^2+w^2)^2} = 0 \]
	
	\[ = (v^2+w^2)(G^2 + H^2) = 0 \]
	
	\[ \Rightarrow G^2 + H^2 = 0 \quad (\because v, w \in \mathbb{R}, \quad v^2 + w^2 > 0)  \]
	
	
	\subsection{Cross terms of $\eta(s)$}
	Now that even for $G$ and $H$, $G^2 + H^2 = 0$ is established, let subscripts $\eta(s)$ denote that such terms are associated to the result of the Dirichlet eta function given $s$ where $s=\frac{1}{b}+ci$.
	
	\[ \eta(s) = \eta(\frac{1}{b}+ci) := G_{\eta(s)}+H_{\eta(s)}i \]
	
	The two self-Cauchy products of $G_{\eta(s)}^2$ and $H_{\eta(s)}^2$ can be expanded and rearranged, similar to how $X_{\zeta(s)}^2$ and $Y_{\zeta(s)}^2$ were processed.
	
	\[ \begin{array}{c}
		G_{\eta(s)}^2 = (g_1 {+} g_2 {+} g_3 {+} \dots {+} g_k)^2 = (g_1^2 {+} g_2^2 {+} g_3^2 {+} \dots {+} g_k^2) \\
		\qquad\qquad\qquad\qquad\qquad\qquad\qquad\qquad\quad {+} 2\{(g_1g_2 {+} g_1g_3 {+} g_1g_4 {+} \dots {+} g_{1}g_{k}) \\ 
		\qquad\qquad\qquad\qquad\qquad\qquad\qquad\qquad {+} (g_2g_3 {+} g_2g_4 {+} g_2g_5 {+} \dots {+} g_2g_{k}) \\ 
		\qquad\qquad\qquad\qquad\qquad\qquad\qquad\qquad {+} (g_3g_4 {+} g_3g_5 {+} g_3g_6 {+} \dots {+} g_3g_{k}) \\ 
		\qquad\qquad\qquad\qquad\qquad\qquad\qquad\qquad \vdots \\
		\qquad\qquad\qquad\qquad {+} \dots {+} (g_{k-1}g_{k}) \} \\ 
	\end{array} \]
	\[ \begin{array}{c}
		\qquad\qquad\qquad\quad = (x_1^2 {+} x_2^2 {+} x_3^2 {+} \dots {+} x_k^2) \\ 
		\qquad\qquad\qquad\qquad\qquad\qquad\quad {+} 2\{(-x_1x_2 {+} x_1x_3 - x_1x_4 {+} \dots \pm x_{1}x_{k}) \\ 
		\qquad\qquad\qquad\qquad\qquad\qquad\qquad {+} (-x_2x_3 {+} x_2x_4 - x_2x_5 {+} \dots \mp x_{2}x_{k}) \\ 
		\qquad\qquad\qquad\qquad\qquad\qquad\qquad {+} (-x_3x_4 {+} x_3x_5 - x_3x_6 {+} \dots \pm x_{3}x_{k}) \\ 
		\qquad\qquad\qquad\qquad\qquad\qquad\qquad \vdots \\
		\qquad\qquad\qquad {+} \dots - (x_{k-1}x_{k}) \} \\ 
	\end{array} \]
	\[\]
	\[\]
	\[ = \Bigl( \frac{\cos({c\ln \frac{1}{1}})}{\sqrt[b]{1}} - \frac{\cos({c\ln \frac{1}{2}})}{\sqrt[b]{2}} + \frac{\cos({c\ln \frac{1}{3}})}{\sqrt[b]{3}} - \dots \pm \frac{\cos({c\ln \frac{1}{k}})}{\sqrt[b]{k}} \Bigr)^2 \]
	\[\]
	\[\]
	\[ = \Bigl( \frac{1}{1} \Bigr)^{\frac2b}\cos^2({c \ln \frac{1}{1}}) + \Bigl( \frac{1}{2} \Bigr)^{\frac2b}\cos^2({c \ln \frac{1}{2}}) + \Bigl( \frac{1}{3} \Bigr)^{\frac2b}\cos^2({c \ln \frac{1}{3}}) + \dots + \Bigl( \frac{1}{k} \Bigr)^{\frac2b}\cos^2({c \ln \frac{1}{k}}) \]
	\[
	\begin{aligned}
		{+}2\Biggl\{ -\Bigl( \frac{1}{1} {\cdot} \frac{1}{2}& \Bigr)^\frac1b \cos(c \ln \frac{1}{1}) \cos(c \ln \frac{1}{2}) {+} \Bigl( \frac{1}{1} {\cdot} \frac{1}{3} \Bigr)^\frac1b \cos(c \ln \frac{1}{1}) \cos(c \ln \frac{1}{3}) \\ 
		& - \Bigl( \frac{1}{1} {\cdot} \frac{1}{4} \Bigr)^\frac1b \cos(c \ln \frac{1}{1}) \cos(c \ln \frac{1}{4}) {+} \dots \pm \Bigl( \frac{1}{1} {\cdot} \frac{1}{k} \Bigr)^\frac1b \cos(c \ln \frac{1}{1}) \cos(c \ln \frac{1}{k}) \\ 		
	\end{aligned}
	\]
	\[
	\begin{aligned}
		\quad - \Bigl( \frac{1}{2} {\cdot} \frac{1}{3}& \Bigr)^\frac1b \cos(c \ln \frac{1}{2}) \cos(c \ln \frac{1}{3}) {+} \Bigl( \frac{1}{2} {\cdot} \frac{1}{4} \Bigr)^\frac1b \cos(c \ln \frac{1}{2}) \cos(c \ln \frac{1}{4}) \\ 
		& - \Bigl( \frac{1}{2} {\cdot} \frac{1}{5} \Bigr)^\frac1b \cos(c \ln \frac{1}{2}) \cos(c \ln \frac{1}{5}) {+} \dots \mp \Bigl( \frac{1}{2} {\cdot} \frac{1}{k} \Bigr)^\frac1b \cos(c \ln \frac{1}{2}) \cos(c \ln \frac{1}{k}) \\
	\end{aligned}
	\]
	\[
	\begin{aligned}
		\quad - \Bigl( \frac{1}{3} {\cdot} \frac{1}{4}& \Bigr)^\frac1b \cos(c \ln \frac{1}{3}) \cos(c \ln \frac{1}{4}) {+} \Bigl( \frac{1}{3} {\cdot} \frac{1}{5} \Bigr)^\frac1b \cos(c \ln \frac{1}{3}) \cos(c \ln \frac{1}{5}) \\
		& - \Bigl( \frac{1}{3} {\cdot} \frac{1}{6} \Bigr)^\frac1b \cos(c \ln \frac{1}{3}) \cos(c \ln \frac{1}{6}) {+} \dots \pm \Bigl( \frac{1}{3} {\cdot} \frac{1}{k} \Bigr)^\frac1b \cos(c \ln \frac{1}{3}) \cos(c \ln \frac{1}{k}) {+} \\ 		
	\end{aligned}
	\]
	\[
	\begin{aligned}
		\vdots \qquad\qquad \vdots \qquad\qquad \vdots \qquad {+} \dots - \Bigl( \frac{1}{(k-1)k} \Bigr)^\frac1b \cos(c \ln \frac{1}{k-1}) \cos(c \ln \frac{1}{k}) \Biggr\}
	\end{aligned}
	\]
	
	\[\]
	\[\]
	\[ \begin{array}{c}
		H_{\eta(s)}^2 = (h_1 {+} h_2 {+} h_3 {+} \dots {+} h_k)^2 = (h_1^2 {+} h_2^2 {+} h_3^2 {+} \dots {+} h_k^2) \\
		\qquad\qquad\qquad\qquad\qquad\qquad\qquad\qquad\quad {+} 2\{(h_1h_2 {+} h_1h_3 {+} h_1h_4 {+} \dots {+} h_{1}h_{k}) \\ 
		\qquad\qquad\qquad\qquad\qquad\qquad\qquad\qquad {+} (h_2h_3 {+} h_2h_4 {+} h_2h_5 {+} \dots {+} h_2h_{k}) \\ 
		\qquad\qquad\qquad\qquad\qquad\qquad\qquad\qquad {+} (h_3h_4 {+} h_3h_5 {+} h_3h_6 {+} \dots {+} h_3h_{k}) \\ 
		\qquad\qquad\qquad\qquad\qquad\qquad\qquad\qquad \vdots \\
		\qquad\qquad\qquad\qquad {+} \dots {+} (h_{k-1}h_{k}) \} \\ 
	\end{array} \]
	\[ \begin{array}{c}
		\qquad\qquad\qquad\quad = (y_1^2 {+} y_2^2 {+} y_3^2 {+} \dots {+} y_k^2) \\ 
		\qquad\qquad\qquad\qquad\qquad\qquad\quad {+} 2\{(-y_1y_2 {+} y_1y_3 - y_1y_4 {+} \dots \pm y_{1}y_{k}) \\ 
		\qquad\qquad\qquad\qquad\qquad\qquad\qquad {+} (-y_2y_3 {+} y_2y_4 - y_2y_5 {+} \dots \mp y_{2}y_{k}) \\ 
		\qquad\qquad\qquad\qquad\qquad\qquad\qquad {+} (-y_3y_4 {+} y_3y_5 - y_3y_6 {+} \dots \pm y_{3}y_{k}) \\ 
		\qquad\qquad\qquad\qquad\qquad\qquad\qquad \vdots \\
		\qquad\qquad\qquad {+} \dots - (y_{k-1}y_{k}) \} \\
	\end{array} \]
	\[\]
	\[\]
	\[ = \Bigl( \frac{\sin({c\ln \frac{1}{1}})}{\sqrt[b]{1}} - \frac{\sin({c\ln \frac{1}{2}})}{\sqrt[b]{2}} {+} \frac{\sin({c\ln \frac{1}{3}})}{\sqrt[b]{3}} - \dots \pm \frac{\sin({c\ln \frac{1}{k}})}{\sqrt[b]{k}} \Bigr)^2 \]
	\[\]
	\[\]
	\[ = \Bigl( \frac{1}{1} \Bigr)^{\frac2b}\sin^2({c \ln \frac{1}{1}}) {+} \Bigl( \frac{1}{2} \Bigr)^{\frac2b}\sin^2({c \ln \frac{1}{2}}) {+} \Bigl( \frac{1}{3} \Bigr)^{\frac2b}\sin^2({c \ln \frac{1}{3}}) {+} \dots {+} \Bigl( \frac{1}{k} \Bigr)^{\frac2b}\sin^2({c \ln \frac{1}{k}}) \]
	\[
	\begin{aligned}
		{+}2\Biggl\{ -\Bigl( \frac{1}{1} {\cdot} \frac{1}{2}& \Bigr)^\frac1b \sin(c \ln \frac{1}{1}) \sin(c \ln \frac{1}{2}) {+} \Bigl( \frac{1}{1} {\cdot} \frac{1}{3} \Bigr)^\frac1b \sin(c \ln \frac{1}{1}) \sin(c \ln \frac{1}{3}) \\ 
		& - \Bigl( \frac{1}{1} {\cdot} \frac{1}{4} \Bigr)^\frac1b \sin(c \ln \frac{1}{1}) \sin(c \ln \frac{1}{4}) {+} \dots \pm \Bigl( \frac{1}{1} {\cdot} \frac{1}{k} \Bigr)^\frac1b \sin(c \ln \frac{1}{1}) \sin(c \ln \frac{1}{k}) \\ 		
	\end{aligned}
	\]
	\[
	\begin{aligned}
		\quad - \Bigl( \frac{1}{2} {\cdot} \frac{1}{3}& \Bigr)^\frac1b \sin(c \ln \frac{1}{2}) \sin(c \ln \frac{1}{3}) {+} \Bigl( \frac{1}{2} {\cdot} \frac{1}{4} \Bigr)^\frac1b \sin(c \ln \frac{1}{2}) \sin(c \ln \frac{1}{4}) \\ 
		& - \Bigl( \frac{1}{2} {\cdot} \frac{1}{5} \Bigr)^\frac1b \sin(c \ln \frac{1}{2}) \sin(c \ln \frac{1}{5}) {+} \dots \mp \Bigl( \frac{1}{2} {\cdot} \frac{1}{k} \Bigr)^\frac1b \sin(c \ln \frac{1}{2}) \sin(c \ln \frac{1}{k}) \\
	\end{aligned}
	\]
	\[
	\begin{aligned}
		\quad - \Bigl( \frac{1}{3} {\cdot} \frac{1}{4}& \Bigr)^\frac1b \sin(c \ln \frac{1}{3}) \sin(c \ln \frac{1}{4}) {+} \Bigl( \frac{1}{3} {\cdot} \frac{1}{5} \Bigr)^\frac1b \sin(c \ln \frac{1}{3}) \sin(c \ln \frac{1}{5}) \\
		& - \Bigl( \frac{1}{3} {\cdot} \frac{1}{6} \Bigr)^\frac1b \sin(c \ln \frac{1}{3}) \sin(c \ln \frac{1}{6}) {+} \dots \pm \Bigl( \frac{1}{3} {\cdot} \frac{1}{k} \Bigr)^\frac1b \sin(c \ln \frac{1}{3}) \sin(c \ln \frac{1}{k}) {+} \\ 		
	\end{aligned}
	\]
	\[
	\begin{aligned}
		\vdots \qquad\qquad \vdots \qquad\qquad \vdots \qquad {+} \dots - \Bigl( \frac{1}{(k-1)k} \Bigr)^\frac1b \sin(c \ln \frac{1}{k-1}) \sin(c \ln \frac{1}{k}) \Biggr\}
	\end{aligned}
	\]
	\[\]
	
	Let the series of cosine cross terms in $G_{\eta(s)}^2$ be $\alpha_{\eta(s)}$ and that of sine cross terms in $H_{\eta(s)}^2$ be $\beta_{\eta(s)}$.
	
	\[ G_{\eta(s)}^2 = g_1^2 + g_2^2 + g_3^2 + \dots + g_k^2 + 2\alpha_{\eta(s)} \]
	\[ = x_1^2 + x_2^2 + x_3^2 + \dots + x_k^2 + 2\alpha_{\eta(s)} \]
	\[ = \sum\limits_{k=1}^{\infty}x_k^2 + 2\alpha_{\eta(s)} \]
	
	\[ H_{\eta(s)}^2 = h_1^2 + h_2^2 + h_3^2 + \dots + h_k^2 + 2\beta_{\eta(s)} \]
	\[ = y_1^2 + y_2^2 + y_3^2 + \dots + y_k^2 + 2\beta_{\eta(s)} \]
	\[ = \sum\limits_{k=1}^{\infty}y_k^2 + 2\beta_{\eta(s)} \]
	
	Adding $G_{\eta(s)}^2$ and $H_{\eta(s)}^2$ together yields a familiar result.
	
	\[ G_{\eta(s)}^2 + H_{\eta(s)}^2 = \sum\limits_{k=1}^{\infty}x_k^2 + 2\alpha_{\eta(s)} + \sum\limits_{k=1}^{\infty}y_k^2 + 2\beta_{\eta(s)} \]
	\[ G_{\eta(s)}^2 + H_{\eta(s)}^2 = \sum\limits_{k=1}^{\infty}(x_k^2 + y_k^2) + 2(\alpha_{\eta(s)} + \beta_{\eta(s)}) \]
	\[ = \sum\limits_{k=1}^{\infty}\frac{1}{k^{\frac{2}{b}}} + 2(\alpha_{\eta(s)} + \beta_{\eta(s)}) \]
	\[ = \zeta(\frac{2}{b}) + 2(\alpha_{\eta(s)} + \beta_{\eta(s)}) = 0 \]
	\[ \zeta(\frac{2}{b}) = - 2(\alpha_{\eta(s)} + \beta_{\eta(s)}) \]
	
	Since $\alpha_{\eta(s)} + \beta_{\eta(s)}$ is the analytically continued version of $\alpha_{\zeta(s)} + \beta_{\zeta(s)}$, $\zeta(s)$ may now be expressed as $\frac{\eta( \frac{2}{b} )}{D_{2/b}}$ where $D_{2/b} = 1 - 2^{1-{2/b}}$.
	\[ \zeta(\frac{2}{b}) = \frac{\eta( \frac{2}{b} )}{D_{2/b}} \]
	
	\section{Convergence and Divergence of $\alpha_{\eta(s)} + \beta_{\eta(s)}$}
	\subsection{Comparing $\alpha_{\eta(s)} + \beta_{\eta(s)}$ to $\alpha_{\zeta(s)} + \beta_{\zeta(s)}$}
	Let us compare $\alpha_{\eta(s)} + \beta_{\eta(s)}$ and $\alpha_{\zeta(s)} + \beta_{\zeta(s)}$.
	
	\[ \alpha_{\zeta(s)} = x_1x_2 + x_1x_3 + x_1x_4 + \dots + x_{1}x_{k} +x_2x_3 + x_2x_4 + x_2x_5 + \dots + x_{2}x_{k} \]
	\[ +x_3x_4 + x_3x_5 + x_3x_6 + \dots + x_{3}x_{k} +x_4x_5 + x_4x_6 + x_4x_7 + \dots + x_{4}x_{k} + \dots + x_{k+1}x_{k} \]
	
	\[ \alpha_{\eta(s)} = -x_1x_2 + x_1x_3 - x_1x_4 + \dots \pm x_{1}x_{k} -x_2x_3 + x_2x_4 - x_2x_5 + \dots \mp x_{2}x_{k} \]
	\[ -x_3x_4 + x_3x_5 - x_3x_6 + \dots \pm x_{3}x_{k} -x_4x_5 + x_4x_6 - x_4x_7 + \dots \mp x_{4}x_{k} - \dots - x_{k-1}x_{k} \]
	
	\[ \beta_{\zeta(s)} = y_1y_2 + y_1y_3 + y_1y_4 + \dots + y_{1}y_{k} +y_2y_3 + y_2y_4 + y_2y_5 + \dots + y_{2}y_{k} \]
	\[ +y_3y_4 + y_3y_5 + y_3y_6 + \dots + y_{3}y_{k} +y_4y_5 + y_4y_6 + y_4y_7 + \dots + y_{4}y_{k} + \dots + y_{k+1}y_{k} \]
	
	\[ \beta_{\eta(s)} = -y_1y_2 + y_1y_3 - y_1y_4 + \dots \pm y_{1}y_{k} -y_2y_3 + y_2y_4 - y_2y_5 + \dots \mp y_{2}y_{k} \]
	\[ -y_3y_4 + y_3y_5 - y_3y_6 + \dots \pm y_{3}y_{k} -y_4y_5 + y_4y_6 - y_4y_7 + \dots \mp y_{4}y_{k} - \dots - y_{k-1}y_{k} \]
	\[\]
	
	$\alpha_{\eta(s)}$ and $\beta_{\eta(s)}$ are of the same sequences as $\alpha_{\zeta(s)}$ and $\beta_{\zeta(s)}$, respectively. However, every cross term with mixed odd-even or even-odd indices has its polarity reversed to a minus. As a result, $\alpha_{\eta(s)} + \beta_{\eta(s)}$ may be rearranged in a similar pattern as $\alpha_{\zeta(s)} + \beta_{\zeta(s)}$ but with alternating polarities for every branch while the signs within every branch alternate per element, always starting with the negative.
	
	\[ \alpha_{\eta(s)} + \beta_{\eta(s)} = \qquad\qquad\qquad\qquad\qquad\qquad\qquad\qquad\qquad\qquad\qquad\qquad (k \to \infty) \]
	\[
	\begin{aligned}
		= \quad \Bigl( \frac{1}{1} \Bigr)^\frac1b \Biggl\{ -\Bigl( \frac{1}{2} \Bigr)^\frac1b & \Bigl( \cos(c \ln \frac{1}{2}) \Bigr) + \Bigl( \frac{1}{3} \Bigr)^\frac1b \Bigl( \cos(c \ln \frac{1}{3}) \Bigr) \\
		& -\Bigl( \frac{1}{4} \Bigr)^\frac1b \Bigl( \cos(c \ln \frac{1}{4}) \Bigr) + \dots \pm \Bigl( \frac{1}{k} \Bigr)^\frac1b \Bigl( \cos(c \ln \frac{1}{k}) \Bigr) \Biggr\} 
	\end{aligned}
	\]
	\[
	\begin{aligned}
		\quad\quad +\Bigl( \frac{1}{2} \Bigr)^\frac1b \Biggl\{ -\Bigl( \frac{1}{3} \Bigr)^\frac1b & \Bigl( \cos(c \ln \frac{2}{3}) \Bigr) + \Bigl( \frac{1}{4} \Bigr)^\frac1b \Bigl( \cos(c \ln \frac{2}{4}) \Bigr) \\
		& -\Bigl( \frac{1}{5} \Bigr)^\frac1b \Bigl( \cos(c \ln \frac{2}{5}) \Bigr) + \dots \mp \Bigl( \frac{1}{k} \Bigr)^\frac1b \Bigl( \cos(c \ln \frac{2}{k}) \Bigr) \Biggr\}
	\end{aligned}
	\]
	\[
	\begin{aligned}
		\quad\quad +\Bigl( \frac{1}{3} \Bigr)^\frac1b \Biggl\{ -\Bigl( \frac{1}{4} \Bigr)^\frac1b & \Bigl( \cos(c \ln \frac{3}{4}) \Bigr) + \Bigl( \frac{1}{5} \Bigr)^\frac1b \Bigl( \cos(c \ln \frac{3}{5}) \Bigr) \\
		& -\Bigl( \frac{1}{6} \Bigr)^\frac1b \Bigl( \cos(c \ln \frac{3}{6}) \Bigr) + \dots \pm \Bigl( \frac{1}{k} \Bigr)^\frac1b \Bigl( \cos(c \ln \frac{3}{k}) \Bigr) \Biggr\} + 
	\end{aligned}
	\]
	\[
	\begin{aligned}
		\quad\quad +\Bigl( \frac{1}{4} \Bigr)^\frac1b \Biggl\{ -\Bigl( \frac{1}{5} \Bigr)^\frac1b & \Bigl( \cos(c \ln \frac{4}{5}) \Bigr) + \Bigl( \frac{1}{6} \Bigr)^\frac1b \Bigl( \cos(c \ln \frac{4}{6}) \Bigr) \\
		& -\Bigl( \frac{1}{7} \Bigr)^\frac1b \Bigl( \cos(c \ln \frac{4}{7}) \Bigr) + \dots \mp \Bigl( \frac{1}{k} \Bigr)^\frac1b \Bigl( \cos(c \ln \frac{4}{k}) \Bigr) \Biggr\} + 
	\end{aligned}
	\]
	\[ \dots + \Bigl( \frac{1}{k-1} \Bigr)^\frac1b \Biggl\{ -\Bigl( \frac{1}{k} \Bigr)^\frac1b \cos(c \ln \frac{k-1}{k}) \Biggr\} \]
	
	
	\subsection{Evaluating the inner Dirichlet series}
	$\alpha_{\eta(s)} + \beta_{\eta(s)}$ is also a nested Dirichlet series. The nodes follow the same pattern from the Riemann zeta function in the form of $\frac{1}{m^{\frac{1}{b}}}$ similar to the nodes in $\alpha_{\zeta(s)} + \beta_{\zeta(s)}$. This means that whether or not $\alpha_{\eta(s)} + \beta_{\eta(s)}$ converges depends on the convergences of branches, the inner Dirichlet series. Every branch is part of a full Dirichlet series with alternating polarity for the summands. With $0 < \frac{1}{b} < 1$ from the initial conditions and along with the newly supposed constraint from the previous iteration of analytic continuation, where $0 < \frac{2}{b} < 1$, every branch will converge. We will denote every branch as $A_m$ and every full Dirichlet series that a branch originates from as $F_m$, where $m \in \mathbb{N}$ and $m \to \infty$, in order to approximate $A_m$.
	
	\[ \alpha_{\eta(s)} + \beta_{\eta(s)} = \Big(\frac{1}{1}\Big)^{\frac{1}{b}}A_1 + \Big(\frac{1}{2}\Big)^{\frac{1}{b}}A_2 + \Big(\frac{1}{3}\Big)^{\frac{1}{b}}A_3 + \dots + \Big(\frac{1}{m}\Big)^{\frac{1}{b}}A_m \]
	
	% \[ F_m = \sum\limits_{k=1}^{\infty} \frac{(-1)^{k-1} {\cdot} \cos(c \ln \frac{m}{k})}{k^{\frac{1}{b}}} {\cdot} (-1)^{m-1} \]
	\[ F_m = \Bigl( \sum\limits_{k=1}^{\infty} \frac{(-1)^{k-1} {\cdot} \cos(c \ln \frac{m}{k})}{k^{\frac{1}{b}}}  \Bigr) {\cdot} (-1)^{m-1} \]
	\[ A_m = \Bigl( \sum\limits_{k=m+1}^{\infty} \frac{(-1)^{k-1} {\cdot} \cos(c \ln \frac{m}{k})}{k^{\frac{1}{b}}} \Bigr) {\cdot} (-1)^{m-1} \]
	\[ = F_m - \Bigl( \sum\limits_{k=1}^{m} \frac{(-1)^{k-1} {\cdot} \cos(c \ln \frac{m}{k})}{k^{\frac{1}{b}}} \Bigr) {\cdot} (-1)^{m-1} \]
	\[\]
	
	$F_m$ and $A_m$ can be expressed in Dirichlet eta function.
	
	\[ \cos(c \ln \frac{m}{k}) = \frac{e^{i \cdot c \ln \frac{m}{k}} + e^{-i \cdot c \ln \frac{m}{k}}}{2} \]
	
	\[ = \frac{e^{\ln (\frac{m}{k})^{ci}} + e^{\ln (\frac{m}{k})^{-ci}}}{2} = \frac{(\frac{m}{k})^{ci} + (\frac{m}{k})^{-ci}}{2} \]
	
	\[ = \frac{m^{ci}k^{-ci} + m^{-ci}k^{ci}}{2} \]
	
	\[ \frac{\cos(c \ln \frac{m}{k}) \cdot (-1)^{k-1}}{k^{\frac{1}{b}}} = \frac{ (-1)^{k-1} {\cdot} (m^{ci}k^{-ci} + m^{-ci}k^{ci}) }{2k^{\frac{1}{b}}}  \]
	
	\[ = \frac{1}{2} \Bigl\{ \frac{ (-1)^{k-1}m^{ci}k^{-ci}  }{k^{\frac{1}{b}}} + \frac{ (-1)^{k-1}m^{-ci}k^{ci}  }{k^{\frac{1}{b}}}  \Bigr\} = \frac{m^{ci}}{2} {\cdot} \frac{(-1)^{k-1}}{k^{\frac{1}{b}+ci}} + \frac{m^{-ci}}{2} {\cdot} \frac{(-1)^{k-1}}{k^{\frac{1}{b}-ci}} \]
	
	%\[ = \frac{1}{2} \Big( m^{ci} {\cdot} \frac{(-1)^{k-1}}{k^{s}} + m^{-ci} {\cdot} \frac{(-1)^{k-1}}{k^{\bar{s}}}   \Big) \]
	\begin{equation} = \frac{1}{2} \Big( m^{ci} {\cdot} \frac{(-1)^{k-1}}{k^{s}} + m^{-ci} {\cdot} \frac{(-1)^{k-1}}{k^{\bar{s}}} \label{eq:1}  \Big) \end{equation}
	
	\[\]
	\[ m^{ci} =  e^{(\ln m){\cdot}ci} = e^{i{\cdot}(c\ln m)} = \cos(c \ln m) + \sin(c \ln m) i \]
	
	\[ m^{-ci} =  e^{(\ln m){\cdot}(-ci)} = e^{i{\cdot}(-c\ln m)} = \cos(-c \ln m) + \sin(-c \ln m) i \]
	
	\[ = \cos(c \ln m) - \sin(c \ln m) i  \]
	\[\]
	
	Each term in $\eta(s)$ and $\eta(\bar{s})$ may be expressed in terms of cosine and sine.
	
	\[ \frac{(-1)^{k-1}}{k^s} = \frac{\cos(c \ln \frac{1}{k}){\cdot}(-1)^{k-1}}{\sqrt[b]{k}} + \frac{\sin(c \ln \frac{1}{k}){\cdot}(-1)^{k-1}}{\sqrt[b]{k}} i \]
	
	\[ \frac{(-1)^{k-1}}{k^{\bar{s}}} = \frac{\cos(-c \ln \frac{1}{k}){\cdot}(-1)^{k-1}}{\sqrt[b]{k}} + \frac{\sin(-c \ln \frac{1}{k}){\cdot}(-1)^{k-1}}{\sqrt[b]{k}} i \]
	
	\[ = \frac{\cos(c \ln \frac{1}{k}){\cdot}(-1)^{k-1}}{\sqrt[b]{k}} - \frac{\sin(c \ln \frac{1}{k}){\cdot}(-1)^{k-1}}{\sqrt[b]{k}} i \]
	\[\]
	
	Replace $m^{ci}$, $m^{-ci}$, $\frac{(-1)^{k-1}}{k^s}$ and $\frac{(-1)^{k-1}}{k^{\bar{s}}}$ in $\eqref{eq:1}$.
	
	\[ m^{ci} {\cdot} \frac{(-1)^{k-1}}{k^s} = \frac{\Big( \cos(c \ln m) {+} \sin(c \ln m) i \Big) \Big(\cos(c \ln \frac{1}{k}) {+} \sin(c \ln \frac{1}{k})i \Big) {\cdot} (-1)^{k-1} }{\sqrt[b]{k}} \]
	
	\[  = \frac{ \Big\{ \cos(c \ln m)\cos(c \ln \frac{1}{k}) {+} \cos(c \ln m)\sin(c \ln \frac{1}{k})i }{\sqrt[b]{k}} \qquad\qquad\qquad\qquad\qquad\qquad\qquad  \]
	\[ \qquad\qquad\qquad\qquad \frac{ {+} \sin(c \ln m)\cos(c \ln \frac{1}{k})i {+} \sin(c \ln m)\sin(c \ln \frac{1}{k})i^2 \Big\} {\cdot} (-1)^{k-1} }{\sqrt[b]{k}}   \]
	
	
	\[ = \frac{ \Big\{ \Big( \cos(c \ln m)\cos(c \ln \frac{1}{k}) - \sin(c \ln m)\sin(c \ln \frac{1}{k}) \Big) }{\sqrt[b]{k}} \qquad\qquad\qquad\qquad\qquad\qquad\qquad  \]
	
	\[ \qquad\qquad\qquad\qquad \frac{ {+} \Big( \cos(c \ln m)\sin(c \ln \frac{1}{k}) {+} \sin(c \ln m)\cos(c \ln \frac{1}{k}) \Big)i \Big\} {\cdot} (-1)^{k-1} }{\sqrt[b]{k}}   \]
	
	\[\]
	\[ m^{-ci} {\cdot} \frac{(-1)^{k-1}}{k^s} = \frac{\Big( \cos(c \ln m) {-} \sin(c \ln m) i \Big) \Big(\cos(c \ln \frac{1}{k}) {-} \sin(c \ln \frac{1}{k})i \Big) {\cdot} (-1)^{k-1} }{\sqrt[b]{k}} \]
	
	\[  = \frac{ \Big\{ \cos(c \ln m)\cos(c \ln \frac{1}{k}) {-} \cos(c \ln m)\sin(c \ln \frac{1}{k})i }{\sqrt[b]{k}} \qquad\qquad\qquad\qquad\qquad\qquad\qquad  \]
	\[ \qquad\qquad\qquad\qquad \frac{ {-} \sin(c \ln m)\cos(c \ln \frac{1}{k})i {+} \sin(c \ln m)\sin(c \ln \frac{1}{k})i^2 \Big\} {\cdot} (-1)^{k-1} }{\sqrt[b]{k}}   \]
	
	
	\[ = \frac{ \Big\{ \Big( \cos(c \ln m)\cos(c \ln \frac{1}{k}) - \sin(c \ln m)\sin(c \ln \frac{1}{k}) \Big)  }{\sqrt[b]{k}} \qquad\qquad\qquad\qquad\qquad\qquad\qquad  \]
	
	\[ \qquad\qquad\qquad\qquad \frac{ {-}\Big( \cos(c \ln m)\sin(c \ln \frac{1}{k}) {+} \sin(c \ln m)\cos(c \ln \frac{1}{k}) \Big)i \Big\} {\cdot} (-1)^{k-1} }{\sqrt[b]{k}}   \]
	\[\]
	\[ \Rightarrow \frac{\cos(c \ln \frac{m}{k}) \cdot (-1)^{k-1}}{k^{\frac{1}{b}}} = \frac{1}{2} \Big( m^{ci} {\cdot} \frac{(-1)^{k-1}}{k^{s}} + m^{-ci} {\cdot} \frac{(-1)^{k-1}}{k^{\bar{s}}}  \Big) \qquad\qquad\qquad\qquad\qquad\qquad \]
	
	\[ = \frac{(-1)^{k-1}}{2} {\cdot} \Biggl\{ \frac{ \Big( \cos(c \ln m)\cos(c \ln \frac{1}{k}) - \sin(c \ln m)\sin(c \ln \frac{1}{k}) \Big) }{\sqrt[b]{k}} \qquad\qquad\qquad\qquad \]
	
	\[ \qquad\qquad\qquad\qquad\qquad\qquad \frac{ {+} \Big( \cos(c \ln m)\sin(c \ln \frac{1}{k}) {+} \sin(c \ln m)\cos(c \ln \frac{1}{k}) \Big)i  }{\sqrt[b]{k}}   \]
	
	\[ \frac{ {+} \Big( \cos(c \ln m)\cos(c \ln \frac{1}{k}) - \sin(c \ln m)\sin(c \ln \frac{1}{k}) \Big)  }{\sqrt[b]{k}}  \qquad \]
	
	\[ \qquad\qquad\qquad\qquad\qquad\qquad \frac{ {-}\Big( \cos(c \ln m)\sin(c \ln \frac{1}{k}) {+} \sin(c \ln m)\cos(c \ln \frac{1}{k}) \Big)i }{\sqrt[b]{k}} \Biggr\}  \]
	\[\]
	\[ = \frac{(-1)^{k-1}}{2} {\cdot} \Bigg\{ \frac{ 2 \Big( \cos(c \ln m)\cos(c \ln \frac{1}{k}) - \sin(c \ln m)\sin(c \ln \frac{1}{k}) \Big) }{\sqrt[b]{k}} \Biggr\}  \]
	\[\]
	\[ = \frac{ \Big( \cos(c \ln m)\cos(c \ln \frac{1}{k}) - \sin(c \ln m)\sin(c \ln \frac{1}{k}) \Big) {\cdot} (-1)^{k-1} }{\sqrt[b]{k}}      \]
	\[\]
	\[ = Re \Big( m^{ci} \cdot \frac{(-1)^{k-1}}{k^{\frac{1}{b}+ci}} \Big) \]
	\\
	\[ F_m = (-1)^{m-1} \cdot \sum\limits_{k=1}^{\infty} \frac{(-1)^{k-1} {\cdot} \cos(c \ln \frac{m}{k})}{k^{\frac{1}{b}}} = (-1)^{m-1} {\cdot} Re \Big( m^{ci} {\cdot} \sum\limits_{k=1}^{\infty}\frac{(-1)^{k-1}}{k^{\frac{1}{b}+ci}}   \Big) \]
	\[ = (-1)^{m-1} {\cdot} Re \Big( m^{ci} {\cdot} \sum\limits_{k=1}^{\infty}\frac{(-1)^{k-1}}{k^s}  \Big) = (-1)^{m-1} \cdot Re \Big( m^{ci} {\cdot} \eta(s)  \Big)   \]
	\\
	Let $R_m(s)$ be the remainder sum when the first $m$ terms are subtracted from $\eta(s)$ and $\eta_m(s)$ be the partial sum of those terms.
	
	\[ R_m(s) = \eta(s) - \eta_m(s) \]
	
	\[ \eta_m(s) = \sum\limits_{k=1}^{m}\frac{(-1)^{k-1}}{k^s} , \quad R_m(s) = \sum\limits_{k=m+1}^{\infty}\frac{(-1)^{k-1}}{k^s} \]
	
	
	
	\[ A_m = (-1)^{m-1} \cdot \sum\limits_{k=m+1}^{\infty} \frac{(-1)^{k-1} {\cdot} \cos(c \ln \frac{m}{k})}{k^{\frac{1}{b}}} = (-1)^{m-1} {\cdot} Re \Big( m^{ci} {\cdot} \sum\limits_{k=m+1}^{\infty}\frac{(-1)^{k-1}}{k^{\frac{1}{b}+ci}}   \Big) \]
	\[ = (-1)^{m-1} {\cdot} Re \Big( m^{ci} {\cdot} \sum\limits_{k=m+1}^{\infty}\frac{(-1)^{k-1}}{k^s}  \Big) = (-1)^{m-1} \cdot Re \Big( m^{ci} {\cdot} R_m(s)  \Big)   \]
	
	Since $s$ is a non-trivial zero for the Dirichlet eta function as well as for the Riemann zeta function, $F_m = 0$ as $\eta(s)$ was defined to converge to zero.
	
	\[ F_m = (-1)^{m-1} \cdot Re(m^{ci} \cdot 0) = 0 \]
	\[ \Rightarrow F_m = 0 \]
	
	$R_m(s)$, the remainder, may be approximated as it asymptotically approaches $\frac{(-1)^{m}}{(m+1)^s}$ as $m$ becomes sufficiently large enough.
	
	\[ R_m \approx \frac{(-1)^{m}}{(m+1)^s} \]
	
	\[ \Rightarrow A_m = (-1)^{m-1} \cdot Re\Bigg(m^{ci} {\cdot} \Big(\frac{(-1)^{m}}{(m+1)^s}\Big) \Bigg) \]
	\\
	The term $Re\Bigg(m^{ci} {\cdot} \Big(\frac{(-1)^{m}}{(m+1)^s}\Big) \Bigg)$ is the first summand of the original branch $A_m$ as $ Re\Bigg(m^{ci} {\cdot} \frac{(-1)^{k-1}}{k^s} \Bigg) = \frac{\cos(c \ln \frac{m}{k}){\cdot}(-1)^{k-1}}{k^{\frac{1}{b}}} $ where $k = m+1$. 
	
	\[ Re\Bigg(m^{ci} {\cdot} \Big(\frac{(-1)^{(m+1)-1}}{(m+1)^s}\Big) \Bigg) = \frac{\cos(c \ln \frac{m}{m+1}){\cdot}(-1)^{(m+1)-1}}{(m+1)^{\frac{1}{b}}}  \]
	\[= \frac{\cos(c \ln \frac{m}{m+1}){\cdot}(-1)^{m}}{(m+1)^{\frac{1}{b}}} \]
	\\
	We process the polarities accordingly for the final form of $A_m$.
	
	\[ A_m = \Biggl( \frac{\cos(c \ln \frac{m}{m+1}){\cdot}(-1)^{m}}{(m+1)^{\frac{1}{b}}}  \Biggr) \cdot (-1)^{m-1} \]
	
	\[ = \Biggl( \frac{\cos(c \ln \frac{m}{m+1})}{(m+1)^{\frac{1}{b}}}  \Biggr) \cdot (-1)^{2m-1} = \frac{-\cos(c \ln \frac{m}{m+1})}{(m+1)^{\frac{1}{b}}} \]	
	\\
	Now $\alpha_{\eta(s)} + \beta_{\eta(s)}$ may be rearranged to determine its convergence.
	\[ \alpha_{\eta(s)} + \beta_{\eta(s)} = \Big(\frac{1}{1}\Big)^{\frac{1}{b}}A_1 + \Big(\frac{1}{2}\Big)^{\frac{1}{b}}A_2 + \Big(\frac{1}{3}\Big)^{\frac{1}{b}}A_3 + \dots + \Big(\frac{1}{m}\Big)^{\frac{1}{b}}A_m \]
	
	\[ = \frac{-\cos(c \ln \frac{1}{2})}{(1{\cdot}2)^{\frac{1}{b}}} + \frac{-\cos(c \ln \frac{2}{3})}{(2{\cdot}3)^{\frac{1}{b}}} + \frac{-\cos(c \ln \frac{3}{4})}{(3{\cdot}4)^{\frac{1}{b}}} + \dots + \frac{-\cos(c \ln \frac{m}{m+1})}{((m(m+1))^{\frac{1}{b}}} \]
	\\
	
	As $m$ grows larger, $m^2$ from $m(m+1) = m^2+m$ will cause the final Dirichlet series to behave as if the exponent of the denominators is $\frac{2}{b}$. This means that it will converge if $\frac{2}{b} > 1$ and diverge if $\frac{2}{b} \leq 1$. 
	
	\section{Divergence of $\frac{\eta(\frac{2}{b})}{D_{2/b}}$}
	%\subsection{Divergence of $\zeta(\frac{2}{b})$}
	Recall that $\zeta(\frac{2}{b}) = \frac{\eta( \frac{2}{b} )}{D_{2/b}} = -2(\alpha_{\eta(s)} + \beta_{\eta(s)}) $. For the purpose of performing analytic continuation when $\zeta(\frac{2}{b}) = -2(\alpha_{\zeta(s)} + \beta_{\zeta(s)})$ because of divergence given any $b$ in $0 < \frac{1}{b} < 1$ or $1 < b < \infty$, we assumed that $\frac{2}{b} \leq 1$ in addition to the initially defined conditions, yielding a new range of $b$ where $b \geq 2$. When given this new range, $\alpha_{\eta(s)} + \beta_{\eta(s)}$ diverges as well as $l {\cdot} (\alpha_{\eta(s)} + \beta_{\eta(s)})$ where $l \in \mathbb{R}_{\neq 0}$.
	\\
	
	In order to find how the newly established relationship between  $\zeta(\frac{2}{b}) = \frac{\eta( \frac{2}{b} )}{D_{2/b}}$ and $-2(\alpha_{\eta(s)} + \beta_{\eta(s)})$ holds true, we look at $\eta(\frac{2}{b})$ and $D_{2/b}$. For $0 < \frac{2}{b} \leq 1$, $\eta(\frac{2}{b})$ always converges. This means that $D_{2/b}$ plays a crucial part in matching the divergence of $-2(\alpha_{\eta(s)} + \beta_{\eta(s)})$. The only way that it can make a convergent value or series diverge is if $D_{2/b} {=} \lim\limits_{n \to \infty}D(\frac{2}{b}) {=} 0$. Below, we may see that zero is within the possible range of $D_{2/b}$.
	
	\[ 0 < \frac{2}{b} < 1 \]
	
	\[ 1 < 2^{\frac{2}{b}} \leq 2 \]
	
	\[ 2^{1-2/b} = 2 / 2^{\frac{2}{b}} \]
	
	\[ 1 \leq 2^{1-2/b} < 2 \]
	
	\[ -2 < -2^{1-2/b} \leq -1 \]
	
	\[ 1-2 < 1-2^{1-2/b} \leq 1-1 \quad \Rightarrow \quad -1 < 1-2^{1-2/b} \leq 0  \] 
	
	\[\Rightarrow -1 < D_{2/b} \leq 0  \] 
	\\
	
	Let every term in $\frac{\eta( \frac{2}{b} )}{D_{\frac{2}{b}}} = - 2(\alpha_{\eta(s)} + \beta_{\eta(s)})$ be represented by either one of the two words: $converge$ or $diverge$.
	\[\]
	\[ \frac{\eta( \frac{2}{b} )}{D_{\frac{2}{b}}} = - 2(\alpha_{\eta(s)} + \beta_{\eta(s)}) \Rightarrow \frac{converge}{D_{\frac{2}{b}}} = - 2 \cdot diverge \]
	
	\[ = \frac{converge}{\lim\limits_{n \to \infty}D(\frac{2}{b})} = - 2 \cdot diverge \]
	
	
	\[ \Rightarrow D_{2/b} = \lim\limits_{n \to \infty}D \Bigl( \frac{2}{b} \Bigr) = 0 \]
	
	\[ \Rightarrow 1-2^{1-2/b} = 0 \]
	
	\[ \Rightarrow 1 = 2^{1-2/b} \]
	
	\[ \Rightarrow 2^0 = 2^{1-2/b} \]
	
	\[ \Rightarrow 0 = 1-2/b \]
	
	\[ \therefore b=2 \ ,  \quad \sigma = \frac{1}{b} = \frac{1}{2} = Re(s) \]
	\\
	$\Rightarrow$ The real part of s where the Riemann zeta function converges to a non-trivial zero is always $\frac{1}{2}$. $\qed$
	
	
\end{document}
